\documentclass[a4paper,reqno]{amsart}

\usepackage[utf8]{inputenc}
\usepackage[english]{babel}
\usepackage{amsmath,amssymb}
\usepackage{hyperref}

\newcommand{\bb}[1]{\lvert #1 \rvert}           % length of a word
\newcommand{\vs}[1]{\mathbin{\downarrow}#1}      % variables of a word
\newcommand{\sig}[1]{\mathbf{#1}}                % signature
\newcommand{\calg}[1]{\boldsymbol{\mathsf{#1}}}  % categorial algebra
\newcommand{\spec}[1]{\boldsymbol{\mathcal{#1}}} % specification
\newcommand{\eeq}{\overset{\mathrm e}{=}}        % existence equation
\newcommand{\QE}{\mathrm{QE}}
\newcommand{\Dom}{\mathrm{Dom}}
\newcommand{\Op}{\mathrm{Op}}
\newcommand{\Eeqt}{\mathrm{Eeqt}}
\newcommand{\id}{\mathrm{id}}
\newcommand{\Id}{\mathrm{Id}}
\newcommand{\gd}{\mathfrak{gd}}
\newcommand{\gl}{\mathrm{gl}}
\newcommand{\src}{\mathrm{sc}}
\newcommand{\tgt}{\mathrm{tg}}
\newcommand{\ii}{\mathrm{i}}
\newcommand{\Gd}{\mathrm{Gd}}
\newcommand{\Uf}{\mathrm{Uf}}
\newcommand{\Ds}{\mathrm{Ds}}
\newcommand{\Tr}{\mathrm{Tr}}
\newcommand{\bGd}{\mathbf{Gd}}
\newcommand{\bUf}{\mathbf{Uf}}
\newcommand{\bDs}{\mathbf{Ds}}
\newcommand{\bTr}{\mathbf{Tr}}
\newcommand{\EPAlg}{\mathsf{EPAlg}}
\newcommand{\PAlg}{\mathsf{PAlg}}
\newcommand{\mor}{\to}
\newcommand{\dmor}{\rightharpoonup}
\newcommand{\T}{\mathrm{T}}
\newcommand{\defi}{\mathrm{def}}

\allowdisplaybreaks
\setlength{\emergencystretch}{3em}

\theoremstyle{plain}
\newtheorem{theorem}{Theorem}[section]
\newtheorem{proposition}[theorem]{Proposition}
\newtheorem{corollary}[theorem]{Corollary}
\newtheorem{lemma}[theorem]{Lemma}
\newtheorem{remark}[theorem]{Remark}

\theoremstyle{definition}
\newtheorem{definition}[theorem]{Definition}
\newtheorem*{convention}{Convention}

\begin{document}

\title{Higher-order categorial algebras}

\author[Climent]{J. Climent Vidal}
\address{Universitat de Val\`{e}ncia\\
         Departament de L\`{o}gica i Filosofia de la Ci\`{e}ncia}
\email{Juan.B.Climent@uv.es}

\author[Cosme]{E. Cosme Ll\'{o}pez}
\address{Universitat de Val\`{e}ncia\\
         Departament de Matem\`{a}tiques\\ and
         Nantong University\\
         School of Mathematics and Statistics}
\email{Enric.Cosme@uv.es}

\author[Ruiz]{R. Ruiz Mora}
\address{Universitat de Val\`{e}ncia\\
         Departament de Matem\`{a}tiques}
\email{Raul.Ruiz-Mora@uv.es}

\author[Vago]{M. Vago Marzal}

\begin{abstract}
We introduce generalized hep-derivors, morphisms between hierarchical equationally partial signatures that send every sort to a set of existence equations in one variable and every operation symbol to a term.
Every generalized hep-derivor induces a functor between the corresponding categories of partial algebras, and the hierarchical equationally partial signatures and the generalized hep-derivors form a category.
We apply these notions to higher-order categorial algebras, i.e., many-sorted algebras whose underlying many-sorted sets are strict higher-order categories and whose operations are higher-order functors.
We present them, in the single-set and in the globular forms, as ep-varieties, and we prove that, in the single-set form, the finite dimension of the cells in dimension $\omega$ cannot be expressed by quasi-existence equations.
By means of generalized hep-derivors we define the graduation, unification, discretization and truncation functors.
We prove that graduation and unification are equivalences, that discretization is left adjoint to truncation, that the categories of $\omega$-categorial algebras are projective limits, that these functors form algebraic squares and that all of them are compatible with the change of signature.
\end{abstract}

\maketitle

\section{Introduction}\label{SIntroduction}

Higher-order categorial algebras are many-sorted algebras whose underlying many-sorted sets carry a structure of strict higher-order category and whose operations are higher-order functors.
They arise in the study of higher-order many-sorted rewriting systems: the paths of a rewriting system of order $n$, up to suitable congruences, form an $n$-categorial algebra, and the higher-order Curry-Howard isomorphisms of~\cite{CVCLRM26} are isomorphisms between such algebras.

There are two natural ways of presenting them.
In the \emph{single-set} presentation~\cite{Str87}, all the cells of a given sort form a single set.
In the \emph{globular} presentation, see~\cite{abgmmm23}, the cells of every dimension form a separate sort.
In both cases the composition is partial, and both presentations are described by quasi-existence equations in the sense of Burmeister~\cite{pb86, pb02}, over equationally partial signatures in the sense of Reichel~\cite{Rei87}.

Our aim is to relate the presentations and the dimensions by functors induced by morphisms between signatures.
To this end we introduce the generalized hep-derivors, which combine the derivors of~\cite{CVST10b}, which send operation symbols to terms, and the general derivors of~\cite{LG18}, which send sorts to formulas.

\textbf{Contents.}
In Section~\ref{SPreliminaries} we fix the notation.
In Section~\ref{SGenHepDerivors} we introduce the generalized hep-derivors and their induced functors.
In Section~\ref{SCatAlg} we present the higher-order categorial algebras as ep-varieties.
In Section~\ref{SGraduation} we define graduation and unification, and in Section~\ref{STruncation} discretization and truncation.
In Section~\ref{SOmega} we consider the dimension $\omega$, and in Section~\ref{SSquares} we assemble these functors in algebraic squares.
Finally, in Section~\ref{SChangeSig} we prove that all of them are compatible with the change of signature.

\section{Preliminaries}\label{SPreliminaries}

We follow the notation of~\cite{CVST10b, CVCLRM26}, to which we refer for the details.

\textbf{Many-sorted sets.}
Let $S$ be a set of sorts.
An $S$-sorted set is a family $A=(A_{s})_{s\in S}$, and an $S$-sorted mapping $f\colon A\mor B$ is a family $(f_{s}\colon A_{s}\mor B_{s})_{s\in S}$.
For a mapping $\varphi\colon T\mor S$, we let $A_{\varphi}$ stand for the $T$-sorted set $(A_{\varphi(t)})_{t\in T}$, and similarly for $f_{\varphi}$.
We let $S^{\star}$ stand for the set of the words on $S$, $\lambda$ for the empty word and $\bb{\mathbf{s}}$ for the length of $\mathbf{s}$.
For $\mathbf{s}\in S^{\star}$, $A_{\mathbf{s}}=\prod_{j\in\bb{\mathbf{s}}}A_{s_{j}}$.
For $\varphi\colon S\mor T$, $\varphi^{\star}\colon S^{\star}\mor T^{\star}$ acts letterwise.

\textbf{Signatures and terms.}
A \emph{many-sorted signature} is a pair $\sig{\Sigma}=(S,\Sigma)$, where $\Sigma=(\Sigma_{\mathbf{s},s})_{(\mathbf{s},s)\in S^{\star}\times S}$ is an $S^{\star}\times S$-sorted set of \emph{operation symbols}.
For an $S$-sorted set $X$ of variables, $\T_{\Sigma}(X)$ is the $S$-sorted set of the terms, and every term is either a variable or $\sigma((P_{j})_{j\in\bb{\mathbf{s}}})$, uniquely, for $\sigma\in\Sigma_{\mathbf{s},s}$ and terms $P_{j}$ of sort $s_{j}$.
Hence we can prove properties of terms by induction and define mappings on terms by recursion.
For $\mathbf{s}\in S^{\star}$, $\vs{\mathbf{s}}$ is the $S$-sorted set of variables $\{v^{\mathbf{s}}_{j}\mid j\in\bb{\mathbf{s}}\}$, where $v^{\mathbf{s}}_{j}$ has sort $s_{j}$.
We identify a family $(a_{j})_{j\in\bb{\mathbf{s}}}\in A_{\mathbf{s}}$ with the valuation of $\vs{\mathbf{s}}$ in $A$ that sends $v^{\mathbf{s}}_{j}$ to $a_{j}$.
For a family of terms $\mathcal{Q}=(Q_{j})_{j\in\bb{\mathbf{s}}}\in\T_{\Sigma}(X)_{\mathbf{s}}$, we let $\mathcal{Q}^{\sharp}$, or $(Q_{j})^{\sharp}_{j\in\bb{\mathbf{s}}}$, stand for the \emph{substitution} of $Q_{j}$ for $v^{\mathbf{s}}_{j}$, i.e., the unique homomorphism from $\T_{\Sigma}(\vs{\mathbf{s}})$ to $\T_{\Sigma}(X)$ that sends $v^{\mathbf{s}}_{j}$ to $Q_{j}$.
The composite of two substitutions is a substitution.
We let $\Op(P)$ stand for the set of the operation symbols that occur in a term $P$, and $\Op(\mathcal{H})$ for those that occur in a set $\mathcal{H}$ of pairs of terms.

\textbf{Derivors.}
A \emph{derivor} from $\sig{\Sigma}=(S,\Sigma)$ to $\sig{\Lambda}=(T,\Lambda)$ is a pair $\mathbf{c}=(\varphi,c)$, where $\varphi\colon S\mor T$ and $c$ assigns to every $\sigma\in\Sigma_{\mathbf{s},s}$ a term $c(\sigma)\in\T_{\Lambda}(\vs{\varphi^{\star}(\mathbf{s})})_{\varphi(s)}$, see~\cite{CVST10b}.
It translates terms by the recursion
\begin{enumerate}
\item[(Der1)] $c^{\sharp}(v^{\mathbf{s}}_{j})=v^{\varphi^{\star}(\mathbf{s})}_{j}$; and
\item[(Der2)] $c^{\sharp}(\sigma((R_{j})_{j\in\bb{\mathbf{r}}}))=(c^{\sharp}(R_{j}))^{\sharp}_{j\in\bb{\mathbf{r}}}(c(\sigma))$,
\end{enumerate}
so that $c^{\sharp}$ sends $\T_{\Sigma}(\vs{\mathbf{s}})_{s}$ to $\T_{\Lambda}(\vs{\varphi^{\star}(\mathbf{s})})_{\varphi(s)}$.
For a set $\mathcal{H}$ of pairs of terms, $c^{\sharp}[\mathcal{H}]=\{(c^{\sharp}(P),c^{\sharp}(Q))\mid(P,Q)\in\mathcal{H}\}$.
The composite of $\mathbf{c}$ and $\mathbf{d}=(\psi,d)\colon\sig{\Lambda}\mor\sig{\Omega}$ is $\mathbf{d}\circ\mathbf{c}=(\psi\circ\varphi,d^{\sharp}\circ c)$, and the identity at $\sig{\Sigma}$ is $(\id_{S},\eta^{\Sigma})$, where $\eta^{\Sigma}(\sigma)=\sigma((v^{\mathbf{s}}_{j})_{j\in\bb{\mathbf{s}}})$.

\begin{proposition}\label{PDerivors}
Let $\mathbf{c}\colon\sig{\Sigma}\mor\sig{\Lambda}$, $\mathbf{d}\colon\sig{\Lambda}\mor\sig{\Omega}$ and $\mathbf{e}\colon\sig{\Omega}\mor\sig{\Xi}$ be derivors.
Then the following statements hold:
\begin{enumerate}
\item[(Der3)] $d^{\sharp}(\mathcal{Q}^{\sharp}(P))=d^{\sharp}(\mathcal{Q})^{\sharp}(d^{\sharp}(P))$, for every term $P$ and every family of terms $\mathcal{Q}$; and
\item[(Der4)] $(e^{\sharp}\circ d)^{\sharp}=e^{\sharp}\circ d^{\sharp}$, $(\eta^{\Lambda})^{\sharp}$ is the identity and $c^{\sharp}\circ\eta^{\Sigma}=c$.
\end{enumerate}
Hence the many-sorted signatures and the derivors form a category $\mathsf{Sig}_{\mathfrak{d}}$.
\end{proposition}
\begin{proof}
We prove (Der3) by induction on $P$.
For a variable $v_{j}$, both sides are $d^{\sharp}(Q_{j})$.
For $P=\tau((P_{i})_{i})$, we have that $\mathcal{Q}^{\sharp}(P)=\tau((\mathcal{Q}^{\sharp}(P_{i}))_{i})$.
Then, by (Der2), the induction hypothesis and the fact that substitutions compose, the following chain of equalities holds
$$
d^{\sharp}(\mathcal{Q}^{\sharp}(P))
=\left(d^{\sharp}(\mathcal{Q})^{\sharp}(d^{\sharp}(P_{i}))\right)^{\sharp}_{i}(d(\tau))
=d^{\sharp}(\mathcal{Q})^{\sharp}\left((d^{\sharp}(P_{i}))^{\sharp}_{i}(d(\tau))\right)
=d^{\sharp}(\mathcal{Q})^{\sharp}(d^{\sharp}(P)).
$$

Now we prove (Der4).
The first equality follows by induction on terms, from (Der2) and (Der3).
The other two follow from (Der1) and (Der2).
They state that the composition is associative and that the identities are neutral.

This completes the proof.
\end{proof}

By (Der4), $e^{\sharp}\circ d^{\sharp}$ only depends on the composite derivor.
Hence, to prove an equality $e^{\sharp}\circ d^{\sharp}=e'^{\sharp}\circ d'^{\sharp}$, it suffices to compare the composite derivors on sorts and on operation symbols.

\textbf{Partial algebras.}
A \emph{partial $\Sigma$-algebra} $\mathbf{A}$ is an $S$-sorted set $A$ together with a partial mapping $\sigma^{\mathbf{A}}\colon A_{\mathbf{s}}\dmor A_{s}$ for every $\sigma\in\Sigma_{\mathbf{s},s}$.
A \emph{homomorphism} $h\colon\mathbf{A}\mor\mathbf{B}$ is an $S$-sorted mapping such that, if $\mathbf{a}\in\Dom(\sigma^{\mathbf{A}})$, then $h\circ\mathbf{a}\in\Dom(\sigma^{\mathbf{B}})$ and $h_{s}(\sigma^{\mathbf{A}}(\mathbf{a}))=\sigma^{\mathbf{B}}(h\circ\mathbf{a})$.
They form a category $\PAlg(\Sigma)$; its full subcategory of the total algebras is $\mathsf{Alg}(\Sigma)$.
A homomorphism is an isomorphism if, and only if, it is bijective and reflects definedness.
Every term $P\in\T_{\Sigma}(X)_{s}$ induces a \emph{partial term operation} $P^{\mathbf{A}}$ from $A_{X}$ to $A_{s}$, defined by recursion on $P$: a variable is always defined, and $\sigma((P_{j})_{j})^{\mathbf{A}}(a)$ is defined if, and only if, every $P_{j}^{\mathbf{A}}(a)$ is defined and $(P_{j}^{\mathbf{A}}(a))_{j}\in\Dom(\sigma^{\mathbf{A}})$, and then it is $\sigma^{\mathbf{A}}((P_{j}^{\mathbf{A}}(a))_{j})$.
Homomorphisms preserve the defined partial term operations.

\textbf{Equations.}
An \emph{existence equation} $P\eeq Q$ is a pair of terms of the same sort.
It is \emph{satisfied} by a valuation $a$ in $\mathbf{A}$, in symbols $\mathbf{A}\models(P,Q)[a]$, if $P^{\mathbf{A}}(a)$ and $Q^{\mathbf{A}}(a)$ are defined and equal.
In particular, $P\eeq P$ states that $P$ is defined.
We let $\Eeqt(\Sigma)_{X}$ stand for the $S$-sorted set of the existence equations with variables in $X$.
For a set $\mathcal{H}$ of existence equations, $\mathbf{A}\models\mathcal{H}[a]$ means that $a$ satisfies all of them.
A \emph{quasi-existence equation} $\bigwedge_{j}(P_{j}\eeq Q_{j})\to(P\eeq Q)$ \emph{holds} in $\mathbf{A}$ if every valuation that satisfies the premisses satisfies the conclusion.
The classes of partial algebras defined by quasi-existence equations, the \emph{$\QE$-varieties}, are closed under products, see~\cite{pb86}.
Homomorphisms preserve satisfaction: if $\mathbf{A}\models\mathcal{H}[a]$, then $\mathbf{B}\models\mathcal{H}[h\circ a]$.
For a substitution $\mathcal{Q}^{\sharp}$, $\mathbf{A}\models\mathcal{Q}^{\sharp}[\mathcal{H}][a]$ if, and only if, every $Q_{j}^{\mathbf{A}}(a)$ is defined and $\mathbf{A}\models\mathcal{H}[(Q_{j}^{\mathbf{A}}(a))_{j}]$.

\textbf{Categories.}
For an adjunction $(G,F,\eta,\varepsilon)$, with $F$ left adjoint to $G$, the triangular identities are $(G\ast\varepsilon)\circ(\eta\ast G)=\id_{G}$ and $(\varepsilon\ast F)\circ(F\ast\eta)=\id_{F}$, see~\cite{sM98}.
A functor $F$ \emph{restricts} to full subcategories $\mathsf{C}'\subseteq\mathsf{C}$ and $\mathsf{D}'\subseteq\mathsf{D}$ if it sends the objects of $\mathsf{C}'$ to objects of $\mathsf{D}'$; we denote the restriction by the same symbol.
A \emph{split $\mathsf{I}$-indexed category} is a functor from $\mathsf{I}^{\mathrm{op}}$ to the category of categories.

\section{Generalized hep-derivors}\label{SGenHepDerivors}

In this section we introduce the hierarchical equationally partial signatures, the generalized hep-derivors between them and the functors that they induce.

\subsection{Hierarchical equationally partial signatures}

\begin{definition}\label{DEpSignature}
An \emph{ep-signature} is a triple $\sig{\Sigma}=(S,\Sigma,\defi^{\Sigma})$, where $(S,\Sigma)$ is a many-sorted signature and $\defi^{\Sigma}$ assigns to every $\sigma\in\Sigma_{\mathbf{s},s}$ a set $\defi^{\Sigma}(\sigma)\subseteq\Eeqt(\Sigma)_{\vs{\mathbf{s}}}$, the \emph{domain conditions} of $\sigma$.
A partial $\Sigma$-algebra $\mathbf{A}$ is a \emph{$\sig{\Sigma}$-algebra} if, for every $\sigma\in\Sigma_{\mathbf{s},s}$,
$$
\Dom(\sigma^{\mathbf{A}})=A_{\defi(\sigma)}=\{\mathbf{a}\in A_{\mathbf{s}}\mid\mathbf{A}\models\defi^{\Sigma}(\sigma)[\mathbf{a}]\},
$$
and a \emph{weak $\sig{\Sigma}$-algebra} if $\Dom(\sigma^{\mathbf{A}})\subseteq A_{\defi(\sigma)}$.
We let $\EPAlg(\sig{\Sigma})\subseteq\mathsf{wEPAlg}(\sig{\Sigma})\subseteq\PAlg(\Sigma)$ stand for the corresponding full subcategories.
The ep-signature $\sig{\Sigma}$ is \emph{hierarchical}, a \emph{hep-signature}, if the relation $\tau\lhd\sigma$ defined by $\tau\in\Op(\defi^{\Sigma}(\sigma))$ is well-founded.
\end{definition}

A many-sorted signature $(S,\Sigma)$ is identified with the hep-signature $(S,\Sigma,\varnothing)$, whose algebras are the total $\Sigma$-algebras.
In a hep-signature, properties of the operation symbols can be proved by well-founded induction along $\lhd$: it suffices to prove the property for $\sigma$ assuming it for every $\tau\lhd\sigma$.
No separate basis is needed, since, for the $\lhd$-minimal symbols, the assumption is vacuous.
In the same way, data can be defined by well-founded recursion along $\lhd$.

\begin{definition}\label{DEpSpec}
An \emph{ep-specification} is a pair $(\sig{\Sigma},\mathcal{E})$, written $\spec{E}$, where $\sig{\Sigma}$ is an ep-signature and $\mathcal{E}$ a set of quasi-existence equations of $\Sigma$; it is a \emph{hep-specification} if $\sig{\Sigma}$ is hierarchical.
We let $\EPAlg(\spec{E})$ stand for the full subcategory of $\EPAlg(\sig{\Sigma})$ whose objects are the $\sig{\Sigma}$-algebras in which every equation of $\mathcal{E}$ holds, the \emph{ep-variety} of $\spec{E}$.
\end{definition}

\begin{convention}
We write the hep-signatures as $\sig{\Sigma}=(S,\Sigma,\defi^{\Sigma})$, $\sig{\Lambda}=(T,\Lambda,\defi^{\Lambda})$ and $\sig{\Omega}=(U,\Omega,\defi^{\Omega})$, and the generalized hep-derivors of Definition~\ref{DGenHepDerivor} as $\mathbf{c}=(\varphi,c,\Dom^{\mathbf{c}})$ and $\mathbf{d}=(\psi,d,\Dom^{\mathbf{d}})$.
\end{convention}

\subsection{Generalized hep-derivors and their functors}

\begin{definition}\label{DGenHepDerivor}
A \emph{generalized hep-derivor} from $\sig{\Sigma}$ to $\sig{\Lambda}$ is a triple $\mathbf{c}=(\varphi,c,\Dom^{\mathbf{c}})$, where $(\varphi,c)$ is a derivor from $(S,\Sigma)$ to $(T,\Lambda)$ and $\Dom^{\mathbf{c}}$ assigns to every $s\in S$ a set $\Dom^{\mathbf{c}}_{s}\subseteq\Eeqt(\Lambda)_{\vs{(\varphi(s))}}$, the \emph{domain} of $s$.
For a partial $\Lambda$-algebra $\mathbf{B}$, the \emph{cells of sort $s$} are the elements $b\in B_{\varphi(s)}$ such that $\mathbf{B}\models\Dom^{\mathbf{c}}_{s}[b]$.
We let $B_{\mathbf{c}}$ stand for the $S$-sorted set of the cells.
\end{definition}

Homomorphisms preserve satisfaction.
Hence, for every homomorphism $f\colon\mathbf{B}\mor\mathbf{B}'$, $f_{\varphi}$ sends cells to cells, and we let $f_{\mathbf{c}}\colon B_{\mathbf{c}}\mor B'_{\mathbf{c}}$ stand for its birestriction.
If all the domains are empty, then $B_{\mathbf{c}}=B_{\varphi}$ and $f_{\mathbf{c}}=f_{\varphi}$.

\begin{proposition}\label{PTransformedAlg}
Let $\mathbf{c}$ be a generalized hep-derivor from $\sig{\Sigma}$ to $\sig{\Lambda}$ and $\mathbf{B}$ a partial $\Lambda$-algebra.
Then there exists a unique partial $\Sigma$-algebra $\mathbf{c}^{\ast}(\mathbf{B})$ on $B_{\mathbf{c}}$ in which, for every $\sigma\in\Sigma_{\mathbf{s},s}$, $\mathbf{b}\in(B_{\mathbf{c}})_{\mathbf{s}}$ is in $\Dom(\sigma^{\mathbf{c}^{\ast}(\mathbf{B})})$ if, and only if,
\begin{enumerate}
\item[(D1)] $\mathbf{c}^{\ast}(\mathbf{B})\models\defi^{\Sigma}(\sigma)[\mathbf{b}]$;
\item[(D2)] $\mathbf{b}\in\Dom(c(\sigma)^{\mathbf{B}})$; and
\item[(D3)] $c(\sigma)^{\mathbf{B}}(\mathbf{b})$ is a cell of sort $s$,
\end{enumerate}
and then $\sigma^{\mathbf{c}^{\ast}(\mathbf{B})}(\mathbf{b})=c(\sigma)^{\mathbf{B}}(\mathbf{b})$.
It is a weak $\sig{\Sigma}$-algebra, the \emph{transformed algebra} of $\mathbf{B}$ along $\mathbf{c}$.
\end{proposition}
\begin{proof}
We define $\sigma^{\mathbf{c}^{\ast}(\mathbf{B})}$ by well-founded recursion along $\lhd$.
The symbols that occur in $\defi^{\Sigma}(\sigma)$ are the $\tau\lhd\sigma$.
Hence condition (D1) only depends on the operations $\tau^{\mathbf{c}^{\ast}(\mathbf{B})}$ already defined.
Conditions (D2) and (D3) only depend on $\mathbf{B}$.
By (D3), the values are cells.
Finally, by (D1), the transformed algebra is a weak $\sig{\Sigma}$-algebra.
\end{proof}

\begin{proposition}\label{PTransformedEpAlg}
Let $\mathbf{c}$ be a generalized hep-derivor from $\sig{\Sigma}$ to $\sig{\Lambda}$ and $\mathbf{B}$ a partial $\Lambda$-algebra.
Then $\mathbf{c}^{\ast}(\mathbf{B})$ is a $\sig{\Sigma}$-algebra if, and only if, for every $\sigma\in\Sigma_{\mathbf{s},s}$ and every family of cells $\mathbf{b}$ that satisfies $\defi^{\Sigma}(\sigma)$ in $\mathbf{c}^{\ast}(\mathbf{B})$, $c(\sigma)^{\mathbf{B}}(\mathbf{b})$ is defined and is a cell of sort $s$.
\end{proposition}
\begin{proof}
By Proposition~\ref{PTransformedAlg}, $\Dom(\sigma^{\mathbf{c}^{\ast}(\mathbf{B})})$ is the set of the families of cells that satisfy (D1), (D2) and (D3).
It is equal to the set of the families that satisfy (D1) if, and only if, every such family satisfies (D2) and (D3).
\end{proof}

\begin{proposition}\label{PGenHepDerivorFunctor}
Let $\mathbf{c}$ be a generalized hep-derivor from $\sig{\Sigma}$ to $\sig{\Lambda}$.
Then the assignment $\mathbf{c}^{\ast}$ that sends $\mathbf{B}$ to $\mathbf{c}^{\ast}(\mathbf{B})$ and $f$ to $f_{\mathbf{c}}$ is a functor from $\PAlg(\Lambda)$ to $\mathsf{wEPAlg}(\sig{\Sigma})$.
\end{proposition}
\begin{proof}
Clearly, $f\mapsto f_{\mathbf{c}}$ preserves identities and composites.
It remains to prove that, for a homomorphism $f\colon\mathbf{B}\mor\mathbf{B}'$, $h=f_{\mathbf{c}}$ is a homomorphism.
We prove, by well-founded induction along $\lhd$, that $h$ preserves the defined operation $\sigma$, assuming it for every $\tau\lhd\sigma$.
Let $\mathbf{b}\in\Dom(\sigma^{\mathbf{c}^{\ast}(\mathbf{B})})$.
We check (D1), (D2) and (D3) for $h\circ\mathbf{b}$.
By the induction hypothesis, $h$ preserves the operations that occur in $\defi^{\Sigma}(\sigma)$, so that $h\circ\mathbf{b}$ satisfies $\defi^{\Sigma}(\sigma)$.
Since $f$ is a homomorphism, $c(\sigma)^{\mathbf{B}'}(h\circ\mathbf{b})$ is defined and equal to $f_{\varphi(s)}(c(\sigma)^{\mathbf{B}}(\mathbf{b}))$, which is a cell.
Hence $h\circ\mathbf{b}\in\Dom(\sigma^{\mathbf{c}^{\ast}(\mathbf{B}')})$ and $h_{s}(\sigma^{\mathbf{c}^{\ast}(\mathbf{B})}(\mathbf{b}))=\sigma^{\mathbf{c}^{\ast}(\mathbf{B}')}(h\circ\mathbf{b})$.
\end{proof}

In general, a generalized hep-derivor does not send ep-algebras to ep-algebras.
For a hep-specification $\spec{F}$ of $\sig{\Lambda}$ and a hep-specification $\spec{E}$ of $\sig{\Sigma}$, we will say that $\mathbf{c}^{\ast}$ \emph{restricts} to a functor from $\EPAlg(\spec{F})$ to $\EPAlg(\spec{E})$ if, for every object $\mathbf{B}$ of $\EPAlg(\spec{F})$, $\mathbf{c}^{\ast}(\mathbf{B})$ is a $\sig{\Sigma}$-algebra in which the equations of $\spec{E}$ hold.
In Sections~\ref{SGraduation} and~\ref{STruncation}, we check it by means of Proposition~\ref{PTransformedEpAlg}.

\subsection{The category of hep-signatures}

\begin{definition}\label{DHepComp}
The \emph{composite} of $\mathbf{c}\colon\sig{\Sigma}\mor\sig{\Lambda}$ and $\mathbf{d}\colon\sig{\Lambda}\mor\sig{\Omega}$ is the generalized hep-derivor
$$
\mathbf{d}\circ\mathbf{c}=\left(\psi\circ\varphi,d^{\sharp}\circ c,\Dom^{\mathbf{d}\circ\mathbf{c}}\right),
\quad\text{where}\quad
\Dom^{\mathbf{d}\circ\mathbf{c}}_{s}=\Dom^{\mathbf{d}}_{\varphi(s)}\cup d^{\sharp}\left[\Dom^{\mathbf{c}}_{s}\right].
$$
The \emph{identity} at $\sig{\Sigma}$ is $\id_{\sig{\Sigma}}=(\id_{S},\eta^{\Sigma},\varnothing)$.
\end{definition}

The translation of terms along $\mathbf{d}$ is related to the transformed algebras as follows.

\begin{lemma}\label{LTermComp}
Let $\mathbf{d}\colon\sig{\Lambda}\mor\sig{\Omega}$ be a generalized hep-derivor, $\mathbf{C}$ a partial $\Omega$-algebra, $P\in\T_{\Lambda}(\vs{\mathbf{t}})_{t}$ and $\mathbf{x}$ a family of cells.
If $P^{\mathbf{d}^{\ast}(\mathbf{C})}(\mathbf{x})$ is defined, then $d^{\sharp}(P)^{\mathbf{C}}(\mathbf{x})$ is defined and equal to it.
Hence, if $\mathbf{d}^{\ast}(\mathbf{C})\models\mathcal{H}[\mathbf{x}]$, then $\mathbf{C}\models d^{\sharp}[\mathcal{H}][\mathbf{x}]$.
\end{lemma}
\begin{proof}
By induction on $P$.
For a variable it follows from (Der1).
Let $P=\tau((P_{j})_{j})$ and $\mathbf{y}=(P_{j}^{\mathbf{d}^{\ast}(\mathbf{C})}(\mathbf{x}))_{j}$.
By the induction hypothesis, $\mathbf{y}=(d^{\sharp}(P_{j})^{\mathbf{C}}(\mathbf{x}))_{j}$.
By (Der2) and the behaviour of substitutions, $d^{\sharp}(P)^{\mathbf{C}}(\mathbf{x})=d(\tau)^{\mathbf{C}}(\mathbf{y})$.
By Proposition~\ref{PTransformedAlg}, this is $\tau^{\mathbf{d}^{\ast}(\mathbf{C})}(\mathbf{y})=P^{\mathbf{d}^{\ast}(\mathbf{C})}(\mathbf{x})$.
\end{proof}

\begin{lemma}\label{LSameCells}
Let $\mathbf{c}=(\varphi,c,\Dom)$ and $\mathbf{c}'=(\varphi,c,\Dom')$ be generalized hep-derivors with the same underlying derivor.
If $\Dom$ and $\Dom'$ determine the same cells in $\mathbf{B}$ and in $\mathbf{B}'$, then $\mathbf{c}^{\ast}(\mathbf{B})=\mathbf{c}'^{\ast}(\mathbf{B})$ and $\mathbf{c}^{\ast}(f)=\mathbf{c}'^{\ast}(f)$, for every homomorphism $f\colon\mathbf{B}\mor\mathbf{B}'$.
In particular, $(\id_{\sig{\Sigma}})^{\ast}$ is the identity on weak $\sig{\Sigma}$-algebras.
\end{lemma}
\begin{proof}
By well-founded induction along $\lhd$, both algebras interpret every $\sigma$ in the same way, since conditions (D1), (D2) and (D3) agree.
On homomorphisms, both functors give the same birestriction.
For the identity, the cells are all the elements, $\eta^{\Sigma}(\sigma)^{\mathbf{A}}=\sigma^{\mathbf{A}}$ and, by induction along $\lhd$, the domain of $\sigma$ in $(\id_{\sig{\Sigma}})^{\ast}(\mathbf{A})$ is $A_{\defi(\sigma)}\cap\Dom(\sigma^{\mathbf{A}})=\Dom(\sigma^{\mathbf{A}})$, since $\mathbf{A}$ is weak.
\end{proof}

\begin{proposition}\label{PHepComp}
Let $\mathbf{c}\colon\sig{\Sigma}\mor\sig{\Lambda}$ and $\mathbf{d}\colon\sig{\Lambda}\mor\sig{\Omega}$ be generalized hep-derivors and $\mathbf{C}$ a partial $\Omega$-algebra.
Then the following statements hold:
\begin{enumerate}
\item $\mathbf{c}^{\ast}(\mathbf{d}^{\ast}(\mathbf{C}))$ is a subalgebra of $(\mathbf{d}\circ\mathbf{c})^{\ast}(\mathbf{C})$, i.e., the inclusion $\kappa^{\mathbf{c},\mathbf{d}}_{\mathbf{C}}$ is a homomorphism, natural in $\mathbf{C}$; and
\item if both algebras have the same cells and $\mathbf{c}^{\ast}(\mathbf{d}^{\ast}(\mathbf{C}))$ is a $\sig{\Sigma}$-algebra, then they are equal; if, moreover, $\mathbf{C}'$ satisfies the same hypotheses, then both functors agree on every homomorphism $f\colon\mathbf{C}\mor\mathbf{C}'$.
\end{enumerate}
\end{proposition}
\begin{proof}
(1) We first prove the inclusion of the cells.
Let $x$ be a cell of sort $s$ of $\mathbf{c}^{\ast}(\mathbf{d}^{\ast}(\mathbf{C}))$.
Then $x$ is a cell of $\mathbf{d}^{\ast}(\mathbf{C})$ and $\mathbf{d}^{\ast}(\mathbf{C})\models\Dom^{\mathbf{c}}_{s}[x]$.
By Lemma~\ref{LTermComp}, $\mathbf{C}\models d^{\sharp}[\Dom^{\mathbf{c}}_{s}][x]$, so that $x$ is a cell of $(\mathbf{d}\circ\mathbf{c})^{\ast}(\mathbf{C})$.

Now we prove, by well-founded induction along $\lhd$, that the inclusion preserves the defined operations.
Let $\mathbf{a}\in\Dom(\sigma^{\mathbf{c}^{\ast}(\mathbf{d}^{\ast}(\mathbf{C}))})$.
By the induction hypothesis, $\mathbf{a}$ satisfies (D1) for $\mathbf{d}\circ\mathbf{c}$.
By Lemma~\ref{LTermComp}, $(d^{\sharp}(c(\sigma)))^{\mathbf{C}}(\mathbf{a})$ is defined and equal to $c(\sigma)^{\mathbf{d}^{\ast}(\mathbf{C})}(\mathbf{a})=\sigma^{\mathbf{c}^{\ast}(\mathbf{d}^{\ast}(\mathbf{C}))}(\mathbf{a})$, which, by the inclusion of the cells, is a cell.
Hence $\mathbf{a}$ satisfies (D2) and (D3), and both operations agree on $\mathbf{a}$.
The naturality is clear, since both functors act on homomorphisms by birestriction.

(2) We prove the reverse inclusion of the domains by well-founded induction along $\lhd$.
Let $\mathbf{a}\in\Dom(\sigma^{(\mathbf{d}\circ\mathbf{c})^{\ast}(\mathbf{C})})$.
By (D1) and the induction hypothesis, $\mathbf{a}$ satisfies $\defi^{\Sigma}(\sigma)$ in $\mathbf{c}^{\ast}(\mathbf{d}^{\ast}(\mathbf{C}))$.
Since the latter is a $\sig{\Sigma}$-algebra, $\mathbf{a}$ is in the domain of $\sigma$ there.
By~(1), $\kappa^{\mathbf{c},\mathbf{d}}_{\mathbf{C}}$ is the identity, and the statement on homomorphisms follows from the naturality.
\end{proof}

By Lemma~\ref{LTermComp}, both algebras have the same cells if, and only if, every cell $x$ of $\mathbf{d}^{\ast}(\mathbf{C})$ that satisfies $d^{\sharp}[\Dom^{\mathbf{c}}_{s}]$ in $\mathbf{C}$ satisfies $\Dom^{\mathbf{c}}_{s}$ in $\mathbf{d}^{\ast}(\mathbf{C})$.
This holds trivially if the domains of $\mathbf{c}$ are empty.
We will use Proposition~\ref{PHepComp}~(2) in the following form.

\begin{lemma}\label{LCommSquare}
Let $\mathbf{a}$, $\mathbf{b}$, $\mathbf{a}'$ and $\mathbf{b}'$ be generalized hep-derivors with $\mathbf{b}\circ\mathbf{a}=\mathbf{b}'\circ\mathbf{a}'$, and $\mathbf{B}$ a partial algebra such that $\mathbf{a}^{\ast}(\mathbf{b}^{\ast}(\mathbf{B}))$ and $\mathbf{a}'^{\ast}(\mathbf{b}'^{\ast}(\mathbf{B}))$ are ep-algebras with the same cells as $(\mathbf{b}\circ\mathbf{a})^{\ast}(\mathbf{B})$.
Then $\mathbf{a}^{\ast}(\mathbf{b}^{\ast}(\mathbf{B}))=\mathbf{a}'^{\ast}(\mathbf{b}'^{\ast}(\mathbf{B}))$, and similarly for homomorphisms between such algebras.
\end{lemma}
\begin{proof}
By Proposition~\ref{PHepComp}~(2), both are equal to $(\mathbf{b}\circ\mathbf{a})^{\ast}(\mathbf{B})$.
\end{proof}

\begin{proposition}\label{PHepSigCat}
The hep-signatures and the generalized hep-derivors, with the composition and the identities of Definition~\ref{DHepComp}, form a category $\mathsf{hepSig}_{\gd}$.
\end{proposition}
\begin{proof}
By (Der4), the underlying derivors compose associatively and the identities are neutral.
It remains to compare the domains.
For $\mathbf{e}=(\chi,e,\Dom^{\mathbf{e}})\colon\sig{\Omega}\mor\sig{\Xi}$, since $e^{\sharp}[\,\cdot\,]$ preserves unions and $(e^{\sharp}\circ d)^{\sharp}=e^{\sharp}\circ d^{\sharp}$, both $\Dom^{\mathbf{e}\circ(\mathbf{d}\circ\mathbf{c})}_{s}$ and $\Dom^{(\mathbf{e}\circ\mathbf{d})\circ\mathbf{c}}_{s}$ are equal to
$$
\Dom^{\mathbf{e}}_{\psi(\varphi(s))}\cup e^{\sharp}\left[\Dom^{\mathbf{d}}_{\varphi(s)}\right]\cup e^{\sharp}\left[d^{\sharp}\left[\Dom^{\mathbf{c}}_{s}\right]\right].
$$
Finally, since $(\eta^{\Lambda})^{\sharp}$ is the identity, the domains of $\id_{\sig{\Lambda}}\circ\mathbf{c}$ and $\mathbf{c}\circ\id_{\sig{\Sigma}}$ are those of $\mathbf{c}$.
\end{proof}

\begin{remark}\label{RHepAlgCat}
By Proposition~\ref{PHepComp}, on the weak ep-algebras the induced functors only preserve composites up to the inclusions $\kappa^{\mathbf{c},\mathbf{d}}$.
Hence they form a normal lax $\mathsf{hepSig}_{\gd}$-indexed category, whose Grothendieck construction yields a category $\mathsf{hepAlg}_{\gd}$ of the ep-algebras of all the hep-signatures.
A morphism from $(\sig{\Sigma},\mathbf{A})$ to $(\sig{\Lambda},\mathbf{B})$ is a pair $(\mathbf{c},f)$ with $f\colon\mathbf{A}\mor\mathbf{c}^{\ast}(\mathbf{B})$, and $(\mathbf{d},g)\circ(\mathbf{c},f)=(\mathbf{d}\circ\mathbf{c},\kappa^{\mathbf{c},\mathbf{d}}_{\mathbf{C}}\circ\mathbf{c}^{\ast}(g)\circ f)$.
For a derivor $(\varphi,c)$ between many-sorted signatures, $(\varphi,c,\varnothing)^{\ast}$ sends a total $\Lambda$-algebra $\mathbf{B}$ to the $\Sigma$-algebra on $B_{\varphi}$ in which $\sigma$ acts as $c(\sigma)^{\mathbf{B}}$, i.e., it is the usual functor, see~\cite{CVST10b}.
Thus $\mathsf{hepAlg}_{\gd}$ contains the category $\mathsf{Alg}_{\mathfrak{d}}$ of~\cite{CVST10b}.
\end{remark}

\section{Higher-order categorial algebras}\label{SCatAlg}

In this section we present the higher-order categorial algebras, in their single-set and globular forms, as ep-varieties of hep-specifications.
In both cases, the composition of two cells is defined exactly when a source of one is a target of the other, which is a domain condition.
We write $x\circ y$ for $\circ(x,y)$, and $\mathbf{x}$ for a family $(x_{p})_{p\in\bb{\mathbf{s}}}$ of variables of arity $\mathbf{s}$.
In this section, $\sig{\Sigma}=(S,\Sigma)$ is a many-sorted signature.

\subsection{The single-set presentation}

In the single-set presentation, all the cells of a sort form a single set, and the dimension of a cell is read off from its sources and targets, as in~\cite{Str87}.

\begin{definition}\label{DSigmaCatSingleSet}
Let $n\in\mathbb{N}$.
The \emph{single-set $n$-categorial signature} $\sig{\Sigma}^{(n)}$ is the hep-signature with sorts $S$ whose operation symbols are those of $\Sigma$ and, for every $k\in n$ and every $s\in S$, the \emph{$k$-source} $\src^{k}_{s}$ and the \emph{$k$-target} $\tgt^{k}_{s}$, of arity $(s)$, and the \emph{$k$-composition} $\circ^{k}_{s}$, of arity $(s,s)$, all of them of coarity $s$.
The only nonempty sets of domain conditions are
$$
\defi(\circ^{k}_{s})=\left\{\left(\src^{k}_{s}(v^{s}_{0}),\tgt^{k}_{s}(v^{s}_{1})\right)\right\}.
$$
It is hierarchical, since the domain conditions only involve sources and targets.
\end{definition}

The same clauses, with $\omega$ in place of $n$, define the hep-signature $\sig{\Sigma}^{(\omega)}$, and $\Sigma^{(\omega)}$ is the union of the $\Sigma^{(n)}$.

\begin{definition}\label{DSpecSingleSet}
Let $n\in\mathbb{N}$.
We let $\spec{C}^{(n)}=(\sig{\Sigma}^{(n)},\mathcal{C}^{(n)})$ stand for the hep-specification whose equations are the following ones, for every $s\in S$, every $\sigma\in\Sigma_{\mathbf{s},s}$, every $k\in n$ and every $l<k$ in $n$; we omit the sort subscripts:
\begin{align*}
&\sigma(\mathbf{x})\eeq\sigma(\mathbf{x});
\tag{A0}\\
&\src^{k}(x)\eeq\src^{k}(x),\quad\tgt^{k}(x)\eeq\tgt^{k}(x);
\tag{A$^{k}$1}\\
&\src^{k}(\src^{k}(x))\eeq\tgt^{k}(\src^{k}(x))\eeq\src^{k}(x),\quad\src^{k}(\tgt^{k}(x))\eeq\tgt^{k}(\tgt^{k}(x))\eeq\tgt^{k}(x);
\tag{A$^{k}$2}\\
&x\circ^{k}y\eeq x\circ^{k}y\ \longleftrightarrow\ \src^{k}(x)\eeq\tgt^{k}(y);
\tag{A$^{k}$3}\\
&x\circ^{k}y\eeq x\circ^{k}y\ \longrightarrow\ \src^{k}(x\circ^{k}y)\eeq\src^{k}(y)\wedge\tgt^{k}(x\circ^{k}y)\eeq\tgt^{k}(x);
\tag{A$^{k}$4}\\
&x\circ^{k}\src^{k}(x)\eeq x,\quad\tgt^{k}(x)\circ^{k}x\eeq x;
\tag{A$^{k}$5}\\
&x\circ^{k}y\eeq x\circ^{k}y\wedge y\circ^{k}z\eeq y\circ^{k}z\ \longrightarrow\ (x\circ^{k}y)\circ^{k}z\eeq x\circ^{k}(y\circ^{k}z);
\tag{A$^{k}$6}\\
&\src^{k}(\sigma(\mathbf{x}))\eeq\sigma(\src^{k}(\mathbf{x})),\quad\tgt^{k}(\sigma(\mathbf{x}))\eeq\sigma(\tgt^{k}(\mathbf{x}));
\tag{A$^{k}$7}\\
&\textstyle\bigwedge_{p}x_{p}\circ^{k}y_{p}\eeq x_{p}\circ^{k}y_{p}\ \longrightarrow\ \sigma(\mathbf{x}\circ^{k}\mathbf{y})\eeq\sigma(\mathbf{x})\circ^{k}\sigma(\mathbf{y});
\tag{A$^{k}$8}\\
&\src^{k}(\src^{l}(x))\eeq\src^{l}(\src^{k}(x))\eeq\src^{l}(\tgt^{k}(x))\eeq\src^{l}(x),\ \text{and dually for $\tgt$};
\tag{AB$^{l,k}$1}\\
&x\circ^{l}y\eeq x\circ^{l}y\ \longrightarrow\ \src^{k}(x\circ^{l}y)\eeq\src^{k}(x)\circ^{l}\src^{k}(y)\wedge\tgt^{k}(x\circ^{l}y)\eeq\tgt^{k}(x)\circ^{l}\tgt^{k}(y);
\tag{AB$^{l,k}$2}\\
&x\circ^{l}y\eeq x\circ^{l}y\wedge z\circ^{l}t\eeq z\circ^{l}t\wedge x\circ^{k}z\eeq x\circ^{k}z\wedge y\circ^{k}t\eeq y\circ^{k}t
\\&\qquad\longrightarrow\ (x\circ^{l}y)\circ^{k}(z\circ^{l}t)\eeq(x\circ^{k}z)\circ^{l}(y\circ^{k}t).
\tag{AB$^{l,k}$3}
\end{align*}
Here $\src^{k}(\mathbf{x})=(\src^{k}(x_{p}))_{p}$, $\mathbf{x}\circ^{k}\mathbf{y}=(x_{p}\circ^{k}y_{p})_{p}$, a biconditional stands for two $\QE$-equations and a conjunction in a conclusion for several ones.
\end{definition}

The same equations, for every $k\in\omega$ and every $l<k$ in $\omega$, define the hep-specification $\spec{C}^{(\omega)}=(\sig{\Sigma}^{(\omega)},\mathcal{C}^{(\omega)})$.

Every model of $\mathcal{C}^{(n)}$, and of $\mathcal{C}^{(\omega)}$, is an ep-algebra: by (A0) and (A$^{k}$1), the symbols without domain conditions are total, and, by (A$^{k}$3), $\circ^{k}$ is defined exactly on the pairs that satisfy its domain condition.
Hence the objects of $\EPAlg(\spec{C}^{(n)})$ are the models of $\mathcal{C}^{(n)}$.

\begin{definition}\label{DCatAlg}
Let $n\in\mathbb{N}$.
We let $\mathsf{nCatAlg}(\sig{\Sigma})$ stand for $\EPAlg(\spec{C}^{(n)})$, whose objects are the \emph{$n$-categorial $\sig{\Sigma}$-algebras}.
For $x$ of sort $s$ and $k\in n$, $x$ is a \emph{$k$-cell} if $\src^{k}_{s}(x)=x$; every cell is an \emph{$n$-cell}.
We let $\omega\mathsf{CatAlg}(\sig{\Sigma})$ stand for the full subcategory of $\EPAlg(\spec{C}^{(\omega)})$ whose objects, the \emph{$\omega$-categorial $\sig{\Sigma}$-algebras}, are those in which every cell is a $k$-cell for some $k\in\omega$.
\end{definition}

The axioms (A$^{k}$1)--(A$^{k}$6) state that, for every $k$, $(\src^{k},\tgt^{k},\circ^{k})$ is a single-set category, and the axioms (AB$^{l,k}$1)--(AB$^{l,k}$3) that every two of them form a single-set $2$-category.
Hence every sort carries a single-set $n$-category in the sense of~\cite{Str87}.
By (A0), (A$^{k}$7) and (A$^{k}$8), the operations of $\Sigma$ are total $n$-functors, and the homomorphisms are the $S$-sorted $n$-functors that commute with them.
Thus $\mathsf{nCatAlg}(\sig{\Sigma})$ is the category of many-sorted algebras in $\mathsf{nCat}$; for $S=1$ and $\Sigma=\varnothing$ it is $\mathsf{nCat}$.
For $n=0$, $\Sigma^{(0)}=\Sigma$ and $\mathsf{0CatAlg}(\sig{\Sigma})=\mathsf{Alg}(\Sigma)$.

We collect the elementary facts on cells that we will use.

\begin{lemma}\label{LGradCells}
Let $n\in\mathbb{N}$, $\mathbf{B}$ an object of $\EPAlg(\spec{C}^{(n)})$ and $j<k$ in $n$.
Then the following statements hold:
\begin{enumerate}
\item if $\mathbf{x}$ is a family of $k$-cells, then $\sigma(\mathbf{x})$ is a $k$-cell;
\item $\src^{j}(x)$ and $\tgt^{j}(x)$ are $j$-cells;
\item every $j$-cell is a $k$-cell; and
\item if $x$, $y$ are $k$-cells and $\src^{j}(x)=\tgt^{j}(y)$, then $x\circ^{j}y$ is defined and is a $k$-cell.
\end{enumerate}
\end{lemma}
\begin{proof}
(1) follows from (A$^{k}$7), and (2) from (A$^{j}$2).
For (3), if $u=\src^{j}(u)$, then, by (AB$^{j,k}$1), $\src^{k}(u)=\src^{k}(\src^{j}(u))=\src^{j}(u)=u$.
For (4), $x\circ^{j}y$ is defined by (A$^{j}$3), and, by (AB$^{j,k}$2), $\src^{k}(x\circ^{j}y)=\src^{k}(x)\circ^{j}\src^{k}(y)=x\circ^{j}y$.
\end{proof}

The proof only uses the axioms for the dimensions $j<k$.
Hence the same statements hold, for $j<k$ in $\omega$, in every object of $\EPAlg(\spec{C}^{(\omega)})$.

In dimension $\omega$, the finite dimension of the cells is not equational.

\begin{proposition}\label{PFiniteDimNotQE}
If $S\neq\varnothing$, then there is no set of $\QE$-equations of $\Sigma^{(\omega)}$ whose models are exactly the $\omega$-categorial $\sig{\Sigma}$-algebras.
In particular, $\omega\mathsf{CatAlg}(\sig{\Sigma})$ is a proper full subcategory of $\EPAlg(\spec{C}^{(\omega)})$.
\end{proposition}
\begin{proof}
Since $\QE$-varieties are closed under products, it suffices to find $\omega$-categorial $\sig{\Sigma}$-algebras $\calg{A}_{k}$, for $k\in\omega$, whose product is not $\omega$-categorial.
Let $\mathsf{G}_{k}$ be the single-set $\omega$-category freely generated by a $(k+1)$-cell $g_{k}$, which is not a $k$-cell.
Let $\calg{A}_{k}$ have $\mathsf{G}_{k}$ in every sort and every $\sigma$ interpreted as the constant mapping with value the $0$-cell $\src^{0}(g_{k})$.
This constant mapping is an $\omega$-functor.
Moreover, every cell of $\calg{A}_{k}$ is a $(k+1)$-cell.
In the product, the cell $(g_{k})_{k\in\omega}$ of sort $s\in S$ is not a $j$-cell for any $j$, since $\src^{j}(g_{j})\neq g_{j}$.
Finally, the product is a model of $\mathcal{C}^{(\omega)}$, which consists of $\QE$-equations.
\end{proof}

\subsection{The globular presentation}

In the globular presentation, see~\cite{abgmmm23}, the cells of every dimension form a separate sort.
For $k\in\mathbb{N}$ and $\mathbf{s}\in S^{\star}$, we let $(k,\mathbf{s})$ stand for the word $((k,s_{p}))_{p\in\bb{\mathbf{s}}}$.

\begin{definition}\label{DSigmaCatGlobular}
Let $n\in\mathbb{N}$ and $N=n+1$.
The \emph{globular $n$-categorial signature} $\sig{\Sigma}^{\gl(n)}$ is the hep-signature with sorts $N\times S$ whose operation symbols are, for every $\sigma\in\Sigma_{\mathbf{s},s}$, every $j<k$ in $N$ and every $s\in S$: a copy $\sigma_{k}$ of $\sigma$, of arity $(k,\mathbf{s})$ and coarity $(k,s)$; the \emph{$(j,k)$-source} $\src^{(j,k)}_{s}$ and the \emph{$(j,k)$-target} $\tgt^{(j,k)}_{s}$, of arity $((k,s))$ and coarity $(j,s)$; the \emph{$(k,j)$-identity} $\ii^{(k,j)}_{s}$, of arity $((j,s))$ and coarity $(k,s)$; and the \emph{$j$-composition of $k$-cells} $\circ^{(j,k)}_{s}$, of arity $((k,s),(k,s))$ and coarity $(k,s)$.
The only nonempty sets of domain conditions are
$$
\defi(\circ^{(j,k)}_{s})=\left\{\left(\src^{(j,k)}_{s}(v^{(k,s)}_{0}),\tgt^{(j,k)}_{s}(v^{(k,s)}_{1})\right)\right\}.
$$
\end{definition}

The same clauses, with $N=\omega$, define the hep-signature $\sig{\Sigma}^{\gl(\omega)}$.
For $n\leq m$ in $\mathbb{N}$, $\Sigma^{\gl(n)}$ embeds canonically into $\Sigma^{\gl(m)}$; it also embeds into $\Sigma^{\gl(\omega)}$, which is the union of the images.

\begin{definition}\label{DSpecGlobular}
Let $n\in\mathbb{N}$ and $N=n+1$.
We let $\spec{C}^{\gl(n)}=(\sig{\Sigma}^{\gl(n)},\mathcal{C}^{\gl(n)})$ stand for the hep-specification whose equations are the following ones, for every $\sigma\in\Sigma_{\mathbf{s},s}$, every $h<j<k$ in $N$ and every $s\in S$, with variables $x,y,z,t$ of sort $(k,s)$ and $u,w$ of sort $(j,s)$ or, in (EA$^{h,j,k}$2), of sort $(h,s)$:
\begin{align*}
&\sigma_{k}(\mathbf{x})\eeq\sigma_{k}(\mathbf{x});
\tag{EA0}\\
&\src^{(j,k)}(x)\eeq\src^{(j,k)}(x),\quad\tgt^{(j,k)}(x)\eeq\tgt^{(j,k)}(x),\quad\ii^{(k,j)}(u)\eeq\ii^{(k,j)}(u);
\tag{EA$^{j,k}$1}\\
&\src^{(j,k)}(\ii^{(k,j)}(u))\eeq u,\quad\tgt^{(j,k)}(\ii^{(k,j)}(u))\eeq u;
\tag{EA$^{j,k}$2}\\
&x\circ^{(j,k)}y\eeq x\circ^{(j,k)}y\ \longleftrightarrow\ \src^{(j,k)}(x)\eeq\tgt^{(j,k)}(y);
\tag{EA$^{j,k}$3}\\
&x\circ^{(j,k)}y\eeq x\circ^{(j,k)}y
\\&\qquad\longrightarrow\ \src^{(j,k)}(x\circ^{(j,k)}y)\eeq\src^{(j,k)}(y)\wedge\tgt^{(j,k)}(x\circ^{(j,k)}y)\eeq\tgt^{(j,k)}(x);
\tag{EA$^{j,k}$4}\\
&x\circ^{(j,k)}\ii^{(k,j)}(\src^{(j,k)}(x))\eeq x,\quad\ii^{(k,j)}(\tgt^{(j,k)}(x))\circ^{(j,k)}x\eeq x;
\tag{EA$^{j,k}$5}\\
&x\circ^{(j,k)}y\eeq x\circ^{(j,k)}y\wedge y\circ^{(j,k)}z\eeq y\circ^{(j,k)}z
\\&\qquad\longrightarrow\ (x\circ^{(j,k)}y)\circ^{(j,k)}z\eeq x\circ^{(j,k)}(y\circ^{(j,k)}z);
\tag{EA$^{j,k}$6}\\
&\src^{(j,k)}(\sigma_{k}(\mathbf{x}))\eeq\sigma_{j}(\src^{(j,k)}(\mathbf{x})),\quad\tgt^{(j,k)}(\sigma_{k}(\mathbf{x}))\eeq\sigma_{j}(\tgt^{(j,k)}(\mathbf{x})),
\\&\ii^{(k,j)}(\sigma_{j}(\mathbf{u}))\eeq\sigma_{k}(\ii^{(k,j)}(\mathbf{u}));
\tag{EA$^{j,k}$7}\\
&\textstyle\bigwedge_{p}x_{p}\circ^{(j,k)}y_{p}\eeq x_{p}\circ^{(j,k)}y_{p}
\\&\qquad\longrightarrow\ \sigma_{k}(\mathbf{x}\circ^{(j,k)}\mathbf{y})\eeq\sigma_{k}(\mathbf{x})\circ^{(j,k)}\sigma_{k}(\mathbf{y});
\tag{EA$^{j,k}$8}\\
&\src^{(h,j)}(\src^{(j,k)}(x))\eeq\src^{(h,j)}(\tgt^{(j,k)}(x))\eeq\src^{(h,k)}(x),\ \text{and dually for $\tgt$};
\tag{EA$^{h,j,k}$1}\\
&\ii^{(k,j)}(\ii^{(j,h)}(w))\eeq\ii^{(k,h)}(w);
\tag{EA$^{h,j,k}$2}\\
&x\circ^{(h,k)}y\eeq x\circ^{(h,k)}y
\\&\qquad\longrightarrow\ \src^{(j,k)}(x\circ^{(h,k)}y)\eeq\src^{(j,k)}(x)\circ^{(h,j)}\src^{(j,k)}(y),\ \text{and dually for $\tgt$};
\tag{EA$^{h,j,k}$3}\\
&x\circ^{(j,k)}y\eeq x\circ^{(j,k)}y\wedge z\circ^{(j,k)}t\eeq z\circ^{(j,k)}t
\\&\qquad\wedge x\circ^{(h,k)}z\eeq x\circ^{(h,k)}z\wedge y\circ^{(h,k)}t\eeq y\circ^{(h,k)}t
\\&\qquad\longrightarrow\ (x\circ^{(j,k)}y)\circ^{(h,k)}(z\circ^{(j,k)}t)\eeq(x\circ^{(h,k)}z)\circ^{(j,k)}(y\circ^{(h,k)}t);
\tag{EA$^{h,j,k}$4}\\
&u\circ^{(h,j)}w\eeq u\circ^{(h,j)}w
\\&\qquad\longrightarrow\ \ii^{(k,j)}(u\circ^{(h,j)}w)\eeq\ii^{(k,j)}(u)\circ^{(h,k)}\ii^{(k,j)}(w).
\tag{EA$^{h,j,k}$5}
\end{align*}
\end{definition}

The same equations, for the dimensions in $\omega$, define the hep-specification $\spec{C}^{\gl(\omega)}=(\sig{\Sigma}^{\gl(\omega)},\mathcal{C}^{\gl(\omega)})$.

As in the single-set case, by (EA0), (EA$^{j,k}$1) and (EA$^{j,k}$3), the models of $\mathcal{C}^{\gl(n)}$ are $\sig{\Sigma}^{\gl(n)}$-algebras.

\begin{definition}\label{DGlCatAlg}
For $n\in\mathbb{N}$, we let $\mathsf{nCat}^{\gl}\mathsf{Alg}(\sig{\Sigma})$ stand for $\EPAlg(\spec{C}^{\gl(n)})$, whose objects are the \emph{globular $n$-categorial $\sig{\Sigma}$-algebras}, and we let $\omega\mathsf{Cat}^{\gl}\mathsf{Alg}(\sig{\Sigma})$ stand for $\EPAlg(\spec{C}^{\gl(\omega)})$.
The elements of sort $(k,s)$ are the \emph{$k$-cells} of sort $s$.
\end{definition}

The axioms on $j<k$ and $h<j<k$ are those of the globular $n$-categories of~\cite{abgmmm23}, and (EA0), (EA$^{j,k}$7) and (EA$^{j,k}$8) state that every $\sigma$ is a globular $n$-functor.
In contrast with Proposition~\ref{PFiniteDimNotQE}, in the globular presentation every cell has a dimension, the first coordinate of its sort.

\begin{lemma}\label{LUnifInj}
In every object of $\EPAlg(\spec{C}^{\gl(n)})$, $\src^{(j,k)}_{s}\circ\ii^{(k,j)}_{s}$ and $\tgt^{(j,k)}_{s}\circ\ii^{(k,j)}_{s}$ are identities.
In particular, $\ii^{(k,j)}_{s}$ is injective.
\end{lemma}
\begin{proof}
It is (EA$^{j,k}$2).
\end{proof}

The following fact, which follows from (EA0), (EA$^{j,k}$7) and (EA$^{j,k}$8) by induction on terms, as in Lemma~\ref{LGradCells}~(1), will be used in Section~\ref{SChangeSig}.
For a term $P$ of $\Lambda$ and $k\in N$, we let $P_{k}$ stand for the term of $\Lambda^{\gl(n)}$ obtained by replacing every $\tau$ by $\tau_{k}$ and every variable $v^{\mathbf{t}}_{p}$ by $v^{(k,\mathbf{t})}_{p}$.

\begin{lemma}\label{LTermFunctor}
Let $\sig{\Lambda}=(T,\Lambda)$ be a many-sorted signature and $P$ a term of $\Lambda$.
Then the following statements hold:
\begin{enumerate}
\item in every object of $\EPAlg(\spec{C}^{(n)})$ for $\sig{\Lambda}$, $P^{\mathbf{B}}$ is total and is an $n$-functor, i.e., it commutes with $\src^{k}$ and $\tgt^{k}$ and preserves the defined $k$-compositions; and
\item in every object of $\EPAlg(\spec{C}^{\gl(n)})$ for $\sig{\Lambda}$, the terms $P_{k}$ are total and form a globular $n$-functor, i.e., $\src^{(j,k)}\circ P_{k}=P_{j}\circ\src^{(j,k)}$, $\tgt^{(j,k)}\circ P_{k}=P_{j}\circ\tgt^{(j,k)}$, $\ii^{(k,j)}\circ P_{j}=P_{k}\circ\ii^{(k,j)}$, and $P_{k}$ preserves the defined $j$-compositions.
\end{enumerate}
\end{lemma}
\begin{proof}
(1) By induction on $P$.
For a variable it is trivial.
For $P=\tau((P_{i})_{i})$, by the induction hypothesis and (A0), $P^{\mathbf{B}}$ is total.
Moreover, $\src^{k}(P(\mathbf{x}))=\tau((\src^{k}(P_{i}(\mathbf{x})))_{i})=\tau((P_{i}(\src^{k}(\mathbf{x})))_{i})=P(\src^{k}(\mathbf{x}))$, by (A$^{k}$7) and the induction hypothesis.
If every $x_{p}\circ^{k}y_{p}$ is defined, then, by the induction hypothesis, $P_{i}(\mathbf{x})\circ^{k}P_{i}(\mathbf{y})$ is defined and equal to $P_{i}(\mathbf{x}\circ^{k}\mathbf{y})$; hence, by (A$^{k}$8), $P(\mathbf{x})\circ^{k}P(\mathbf{y})$ is defined and equal to $P(\mathbf{x}\circ^{k}\mathbf{y})$.

(2) The proof is the same, with (EA0), (EA$^{j,k}$7) and (EA$^{j,k}$8), since $(\tau((P_{i})_{i}))_{k}=\tau_{k}(((P_{i})_{k})_{i})$.
\end{proof}

\section{Graduation and unification}\label{SGraduation}

In this section we pass from the single-set to the globular presentation, and back, by means of two generalized hep-derivors.
In what follows, $\sig{\Sigma}=(S,\Sigma)$ is a many-sorted signature and $n\in\mathbb{N}$.
For the generalized hep-derivors of this work we use the following notation: the graduation derivor is $\bGd^{(n)}=(\varphi^{\Gd},\Gd^{(n)},\Dom^{\Gd})$, where $\Gd^{(n)}$ denotes both the assignment of terms and the induced functor, and similarly for the unification, discretization and truncation derivors $\bUf$, $\bDs$ and $\bTr$.
We omit the subscript $\sig{\Sigma}$.

\subsection{Graduation}

Graduating a single-set $n$-categorial algebra consists in distributing its cells into sorts according to their dimension.

\begin{definition}\label{DGradDerivor}
The \emph{graduation derivor} $\bGd^{(n)}$ is the generalized hep-derivor from $\sig{\Sigma}^{\gl(n)}$ to $\sig{\Sigma}^{(n)}$ that sends the sort $(k,s)$ to $s$ and, for every $j<k$ in $n+1$,
\begin{align*}
\Gd^{(n)}(\sigma_{k})&=\sigma((v^{\mathbf{s}}_{p})_{p}), &
\Gd^{(n)}(\ii^{(k,j)}_{s})&=v^{s}_{0}, &
\Gd^{(n)}(\circ^{(j,k)}_{s})&=v^{s}_{0}\circ^{j}_{s}v^{s}_{1},\\
\Gd^{(n)}(\src^{(j,k)}_{s})&=\src^{j}_{s}(v^{s}_{0}), &
\Gd^{(n)}(\tgt^{(j,k)}_{s})&=\tgt^{j}_{s}(v^{s}_{0}),
\end{align*}
with domains $\Dom^{\Gd}_{(k,s)}=\{(v^{s}_{0},\src^{k}_{s}(v^{s}_{0}))\}$, for $k\in n$, and $\Dom^{\Gd}_{(n,s)}=\varnothing$.
The \emph{graduation functor} is $\Gd^{(n)}=(\bGd^{(n)})^{\ast}$.
\end{definition}

Thus the cells of sort $(k,s)$ of $\Gd^{(n)}(\mathbf{B})$ are the $k$-cells of sort $s$ of $\mathbf{B}$, for $k\in n$, and all its cells of sort $s$, for $k=n$.

\begin{proposition}\label{PGradDerivor}
The functor $\Gd^{(n)}$ restricts to a functor from $\mathsf{nCatAlg}(\sig{\Sigma})$ to $\mathsf{nCat}^{\gl}\mathsf{Alg}(\sig{\Sigma})$.
Moreover, in $\Gd^{(n)}(\calg{B})$, $\sigma_{k}$, $\src^{(j,k)}_{s}$ and $\tgt^{(j,k)}_{s}$ act as $\sigma$, $\src^{j}_{s}$ and $\tgt^{j}_{s}$, the identities are inclusions, and $x\circ^{(j,k)}_{s}y$ is defined if, and only if, $x\circ^{j}_{s}y$ is defined, and then it is $x\circ^{j}_{s}y$.
\end{proposition}
\begin{proof}
Let $\calg{B}$ be an $n$-categorial $\sig{\Sigma}$-algebra.
We first prove that $\Gd^{(n)}(\calg{B})$ is a $\sig{\Sigma}^{\gl(n)}$-algebra, by Proposition~\ref{PTransformedEpAlg}.
By (A0) and (A$^{j}$1), the terms assigned to $\sigma_{k}$, $\src^{(j,k)}_{s}$, $\tgt^{(j,k)}_{s}$ and $\ii^{(k,j)}_{s}$ are total.
By Lemma~\ref{LGradCells}, they send cells to cells of the right dimension.
Two $k$-cells satisfy the domain condition of $\circ^{(j,k)}_{s}$ if, and only if, $\src^{j}_{s}(x)=\tgt^{j}_{s}(y)$, and then, by Lemma~\ref{LGradCells}~(4), $x\circ^{j}_{s}y$ is defined and is a $k$-cell.
Hence it is a $\sig{\Sigma}^{\gl(n)}$-algebra, and, by Proposition~\ref{PTransformedAlg}, its operations are as stated.

Now we prove that the axioms of $\mathcal{C}^{\gl(n)}$ hold.
Let $h<j<k$.
The axioms (EA0), (EA$^{j,k}$1) and (EA$^{j,k}$3) hold by the above.
(EA$^{j,k}$2) holds since $\src^{j}(u)=u=\tgt^{j}(u)$ for every $j$-cell $u$.
(EA$^{j,k}$4)--(EA$^{j,k}$8) become (A$^{j}$4)--(A$^{j}$8) restricted to $k$-cells.
(EA$^{h,j,k}$2) and (EA$^{h,j,k}$5) are trivial, since the identities are inclusions.
Finally, (EA$^{h,j,k}$1), (EA$^{h,j,k}$3) and (EA$^{h,j,k}$4) become (AB$^{h,j}$1), (AB$^{h,j}$2) and (AB$^{h,j}$3), after exchanging $y$ and $z$ in the last one.
\end{proof}

The restriction to $\mathsf{nCatAlg}(\sig{\Sigma})$ is necessary: in a $\sig{\Sigma}^{(n)}$-algebra in which $\src^{j}$ is not idempotent, the $j$-source of a cell need not be a $j$-cell, and $\src^{(j,k)}$ need not be total in $\Gd^{(n)}(\mathbf{B})$.

\subsection{Unification}

Unifying a globular $n$-categorial algebra consists in keeping its $n$-cells, in which every $k$-cell is represented by its $(n,k)$-identity.

\begin{definition}\label{DUnifDerivor}
The \emph{unification derivor} $\bUf^{(n)}$ is the generalized hep-derivor from $\sig{\Sigma}^{(n)}$ to $\sig{\Sigma}^{\gl(n)}$, with empty domains, that sends $s$ to $(n,s)$ and, for every $k\in n$,
\begin{align*}
\Uf^{(n)}(\sigma)&=\sigma_{n}((v^{(n,\mathbf{s})}_{p})_{p}), &
\Uf^{(n)}(\src^{k}_{s})&=\ii^{(n,k)}_{s}(\src^{(k,n)}_{s}(v^{(n,s)}_{0})),\\
\Uf^{(n)}(\circ^{k}_{s})&=v^{(n,s)}_{0}\circ^{(k,n)}_{s}v^{(n,s)}_{1}, &
\Uf^{(n)}(\tgt^{k}_{s})&=\ii^{(n,k)}_{s}(\tgt^{(k,n)}_{s}(v^{(n,s)}_{0})).
\end{align*}
The \emph{unification functor} is $\Uf^{(n)}=(\bUf^{(n)})^{\ast}$.
\end{definition}

\begin{proposition}\label{PUnifDerivor}
The functor $\Uf^{(n)}$ restricts to a functor from $\mathsf{nCat}^{\gl}\mathsf{Alg}(\sig{\Sigma})$ to $\mathsf{nCatAlg}(\sig{\Sigma})$.
Moreover, the cells of sort $s$ of $\Uf^{(n)}(\calg{A})$ are the $n$-cells of sort $s$ of $\calg{A}$, in which $\sigma$, $\src^{k}_{s}$, $\tgt^{k}_{s}$ and $\circ^{k}_{s}$ act as $\sigma_{n}$, $\ii^{(n,k)}_{s}\circ\src^{(k,n)}_{s}$, $\ii^{(n,k)}_{s}\circ\tgt^{(k,n)}_{s}$ and $\circ^{(k,n)}_{s}$.
\end{proposition}
\begin{proof}
Let $\calg{A}$ be a globular $n$-categorial $\sig{\Sigma}$-algebra.
We first prove that $\Uf^{(n)}(\calg{A})$ is a $\sig{\Sigma}^{(n)}$-algebra.
The domains are empty, so that the cells are the $n$-cells.
By (EA0) and (EA$^{k,n}$1), the terms assigned to $\sigma$, $\src^{k}_{s}$ and $\tgt^{k}_{s}$ are total.
By Lemma~\ref{LUnifInj}, two $n$-cells satisfy the domain condition of $\circ^{k}_{s}$ if, and only if, $\src^{(k,n)}_{s}(x)=\tgt^{(k,n)}_{s}(y)$, and then, by (EA$^{k,n}$3), $x\circ^{(k,n)}_{s}y$ is defined.
Hence, by Propositions~\ref{PTransformedEpAlg} and~\ref{PTransformedAlg}, it is a $\sig{\Sigma}^{(n)}$-algebra with the stated operations.

Now we prove that the axioms of $\mathcal{C}^{(n)}$ hold.
Let $l<k$ in $n$.
(A0), (A$^{k}$1) and (A$^{k}$3) follow from (EA0), (EA$^{k,n}$1), (EA$^{k,n}$3) and Lemma~\ref{LUnifInj}.
(A$^{k}$2) follows from Lemma~\ref{LUnifInj}.
(A$^{k}$4), (A$^{k}$5), (A$^{k}$6) and (A$^{k}$8) follow from (EA$^{k,n}$4), (EA$^{k,n}$5), (EA$^{k,n}$6) and (EA$^{k,n}$8).
(A$^{k}$7) follows from the first and the third equations of (EA$^{k,n}$7).
By (EA$^{l,k,n}$1) and Lemma~\ref{LUnifInj}, $\src^{(l,n)}(\ii^{(n,k)}(w))=\src^{(l,k)}(w)$ for every $k$-cell $w$; with this, (AB$^{l,k}$1) follows from (EA$^{l,k,n}$1) and (EA$^{l,k,n}$2).
(AB$^{l,k}$2) follows from (EA$^{l,k,n}$3) and (EA$^{l,k,n}$5), since
\begin{align*}
\ii^{(n,k)}(\src^{(k,n)}(x\circ^{(l,n)}y))
&=\ii^{(n,k)}(\src^{(k,n)}(x)\circ^{(l,k)}\src^{(k,n)}(y))
\\&=\ii^{(n,k)}(\src^{(k,n)}(x))\circ^{(l,n)}\ii^{(n,k)}(\src^{(k,n)}(y)).
\end{align*}
Finally, (AB$^{l,k}$3) is (EA$^{l,k,n}$4), after exchanging $y$ and $z$.
\end{proof}

\subsection{Graduation and unification are equivalences}

We compute the composites of the functors by composing the corresponding generalized hep-derivors.

\begin{proposition}\label{PUnifGrad}
The equality $\bGd^{(n)}\circ\bUf^{(n)}=\id_{\sig{\Sigma}^{(n)}}$ holds in $\mathsf{hepSig}_{\gd}$.
Consequently, $\Uf^{(n)}\circ\Gd^{(n)}=\Id_{\mathsf{nCatAlg}(\sig{\Sigma})}$.
\end{proposition}
\begin{proof}
The composite sends $s$ to $s$.
By (Der1) and (Der2), it sends $\src^{k}_{s}$ to $(\Gd^{(n)})^{\sharp}(\ii^{(n,k)}_{s}(\src^{(k,n)}_{s}(v_{0})))=\src^{k}_{s}(v_{0})$, and similarly every other operation symbol to itself, applied to variables.
Its domains are $\Dom^{\Gd}_{(n,s)}\cup(\Gd^{(n)})^{\sharp}[\varnothing]=\varnothing$.
Hence it is the identity.

Now let $\calg{A}$ be an $n$-categorial $\sig{\Sigma}$-algebra.
By Propositions~\ref{PGradDerivor} and~\ref{PUnifDerivor}, the cells of $\Uf^{(n)}(\Gd^{(n)}(\calg{A}))$ are those of $\calg{A}$, and it is a $\sig{\Sigma}^{(n)}$-algebra.
Therefore, by Proposition~\ref{PHepComp}~(2) and Lemma~\ref{LSameCells}, $\Uf^{(n)}(\Gd^{(n)}(\calg{A}))=(\id_{\sig{\Sigma}^{(n)}})^{\ast}(\calg{A})=\calg{A}$, and similarly for homomorphisms.
\end{proof}

The composite in the other order is not the identity, since the $k$-cells of $\Gd^{(n)}(\Uf^{(n)}(\calg{A}))$ are the $(n,k)$-identities of the $k$-cells of $\calg{A}$.

\begin{proposition}\label{PGradUnif}
For every globular $n$-categorial $\sig{\Sigma}$-algebra $\calg{A}$, let $\theta_{\calg{A}}$ send every $k$-cell $u$ of sort $s$ to $\ii^{(n,k)}_{s}(u)$, where $\ii^{(n,n)}_{s}$ is the identity.
Then $\theta$ is a natural isomorphism from $\Id$ to $\Gd^{(n)}\circ\Uf^{(n)}$.
\end{proposition}
\begin{proof}
By Propositions~\ref{PGradDerivor} and~\ref{PUnifDerivor}, the cells of sort $(k,s)$ of $\Gd^{(n)}(\Uf^{(n)}(\calg{A}))$ are the $n$-cells $x$ with $x=\ii^{(n,k)}(\src^{(k,n)}(x))$.
Its $(j,k)$-sources, $(j,k)$-targets and $j$-compositions are $\ii^{(n,j)}\circ\src^{(j,n)}$, $\ii^{(n,j)}\circ\tgt^{(j,n)}$ and $\circ^{(j,n)}$, its identities are inclusions, and $\sigma_{k}$ acts as $\sigma_{n}$.

We first prove that $\theta_{\calg{A}}$ is a bijection onto these cells.
By Lemma~\ref{LUnifInj}, $\ii^{(n,k)}(u)=\ii^{(n,k)}(\src^{(k,n)}(\ii^{(n,k)}(u)))$, so that $\theta_{\calg{A}}(u)$ is a cell of sort $(k,s)$.
Its inverse sends a cell $x$ of sort $(k,s)$ to $\src^{(k,n)}(x)$.

Now we prove that $\theta_{\calg{A}}$ is a homomorphism that reflects definedness.
By the third equation of (EA$^{k,n}$7), it commutes with $\sigma_{k}$.
By (EA$^{j,k,n}$1) and Lemma~\ref{LUnifInj}, $\ii^{(n,j)}(\src^{(j,n)}(\ii^{(n,k)}(u)))=\ii^{(n,j)}(\src^{(j,k)}(u))$, so that it commutes with the sources, and similarly with the targets.
By (EA$^{j,k,n}$2), $\ii^{(n,k)}(\ii^{(k,j)}(w))=\ii^{(n,j)}(w)$, so that it commutes with the identities.
By Lemma~\ref{LUnifInj} and (EA$^{j,k}$3), $\theta(u)\circ^{(j,k)}\theta(u')$ is defined if, and only if, $u\circ^{(j,k)}u'$ is defined, and then, by (EA$^{j,k,n}$5), $\theta(u\circ^{(j,k)}u')=\theta(u)\circ^{(j,n)}\theta(u')$.
Hence $\theta_{\calg{A}}$ is an isomorphism.

Finally, $\theta$ is natural, since the homomorphisms commute with the identities.
\end{proof}

\begin{proposition}\label{PGdUfAdj}
The functors $\Gd^{(n)}$ and $\Uf^{(n)}$ are mutually quasi-inverse equivalences.
More precisely, $(\Uf^{(n)},\Gd^{(n)},\id,\theta^{-1})$ is an adjoint equivalence, with $\Gd^{(n)}$ left adjoint to $\Uf^{(n)}$.
\end{proposition}
\begin{proof}
By Propositions~\ref{PUnifGrad} and~\ref{PGradUnif}, the unit and the counit are natural isomorphisms.
Since the unit is the identity, the triangular identities state that $\Uf^{(n)}(\theta^{-1}_{\calg{A}})$ and $\theta^{-1}_{\Gd^{(n)}(\calg{B})}$ are identities.
The first one is the family of the components of $\theta^{-1}_{\calg{A}}$ at the sorts $(n,s)$, which are identities.
For the second one, the identities of $\Gd^{(n)}(\calg{B})$ are inclusions, so that $\theta_{\Gd^{(n)}(\calg{B})}$ is the identity.
\end{proof}

\section{Discretization and truncation}\label{STruncation}

In this section we move between dimensions.
Let $\sig{\Sigma}=(S,\Sigma)$ be a many-sorted signature and $m\leq n$ in $\mathbb{N}$.
Discretization regards an $m$-categorial algebra as an $n$-categorial algebra whose cells of dimension at least $m$ are identities, and truncation forgets the cells of dimension greater than $m$.

\subsection{The functors}

\begin{definition}\label{DDiscDerivor}
Let $m<n$.
The \emph{discretization derivors} are the following generalized hep-derivors with empty domains:
\begin{enumerate}
\item $\bDs^{(n,m)}$, from $\sig{\Sigma}^{(n)}$ to $\sig{\Sigma}^{(m)}$, with $\id_{S}$ on sorts, which sends $\sigma$, $\src^{k}_{s}$, $\tgt^{k}_{s}$ and $\circ^{k}_{s}$, for $k<m$, to themselves, applied to variables, and $\src^{k}_{s}$, $\tgt^{k}_{s}$ and $\circ^{k}_{s}$, for $k\geq m$, to $v^{s}_{0}$; and
\item $\bDs^{\gl(n,m)}$, from $\sig{\Sigma}^{\gl(n)}$ to $\sig{\Sigma}^{\gl(m)}$, which sends $(k,s)$ to $(\bar{k},s)$, where $\bar{k}=\min\{k,m\}$, $\sigma_{k}$ to $\sigma_{\bar{k}}$, the symbols $\src^{(j,k)}_{s}$, $\tgt^{(j,k)}_{s}$, $\ii^{(k,j)}_{s}$ and $\circ^{(j,k)}_{s}$, for $j<m$, to $\src^{(j,\bar{k})}_{s}$, $\tgt^{(j,\bar{k})}_{s}$, $\ii^{(\bar{k},j)}_{s}$ and $\circ^{(j,\bar{k})}_{s}$, and, for $j\geq m$, to $v^{(m,s)}_{0}$, all of them applied to variables.
\end{enumerate}
For $m=n$, both are identities.
The \emph{discretization functors} are $\Ds^{(n,m)}=(\bDs^{(n,m)})^{\ast}$ and $\Ds^{\gl(n,m)}=(\bDs^{\gl(n,m)})^{\ast}$.
\end{definition}

For $k\geq m$, the term assigned to $\circ^{k}_{s}$ is total.
However, in the transformed algebra $\src^{k}_{s}$ and $\tgt^{k}_{s}$ act as identities, so that, by (D1), $x\circ^{k}_{s}y$ is defined if, and only if, $x=y$.
No term of $\Sigma^{(m)}$ could express this condition.

\begin{proposition}\label{PDiscDerivor}
The functors $\Ds^{(n,m)}$ and $\Ds^{\gl(n,m)}$ restrict to functors from $\mathsf{mCatAlg}(\sig{\Sigma})$ to $\mathsf{nCatAlg}(\sig{\Sigma})$ and from $\mathsf{mCat}^{\gl}\mathsf{Alg}(\sig{\Sigma})$ to $\mathsf{nCat}^{\gl}\mathsf{Alg}(\sig{\Sigma})$.
Moreover, $\Ds^{(n,m)}(\calg{B})$ has the cells and the structure of dimension $k<m$ of $\calg{B}$, and, for $k\geq m$, $\src^{k}$ and $\tgt^{k}$ are identities and $x\circ^{k}y$ is defined if, and only if, $x=y$, and then it is $x$.
The $k$-cells of $\Ds^{\gl(n,m)}(\calg{B})$ are the $\bar{k}$-cells of $\calg{B}$, and its structure is described similarly.
\end{proposition}
\begin{proof}
For $m=n$ it is trivial.
Let $m<n$ and $\calg{B}$ an $m$-categorial $\sig{\Sigma}$-algebra.
We first prove that $\Ds^{(n,m)}(\calg{B})$ is a $\sig{\Sigma}^{(n)}$-algebra.
The domains are empty, so that the cells are all the cells of $\calg{B}$.
By (A0) and (A$^{k}$1), the symbols without domain conditions act as stated.
For $k<m$, by (A$^{k}$3), $x\circ^{k}y$ is defined in $\calg{B}$ when $\src^{k}(x)=\tgt^{k}(y)$.
For $k\geq m$, the domain condition holds if, and only if, $x=y$, and then $v_{0}$ is defined.
Hence, by Proposition~\ref{PTransformedEpAlg}, it is a $\sig{\Sigma}^{(n)}$-algebra.

Now we prove the axioms.
Those in which all the dimensions are smaller than $m$ are axioms of $\mathcal{C}^{(m)}$, which hold in $\calg{B}$.
The remaining ones become trivial, since $\src^{k}$ and $\tgt^{k}$ are identities and $\circ^{k}$ is only defined on equal cells for $k\geq m$.
For instance, if $l<m\leq k$ and $x\circ^{l}y$ is defined, both sides of (AB$^{l,k}$2) are $x\circ^{l}y$.

The globular case is similar.
The axioms in which all the dimensions are at most $m$ are axioms of $\mathcal{C}^{\gl(m)}$, those with $j<m<k$ become axioms with $k$ replaced by $m$, and those with a lower dimension $j\geq m$ become trivial.
\end{proof}

\begin{definition}\label{DTruncDerivor}
Let $m<n$.
The \emph{truncation derivors} are the following generalized hep-derivors:
\begin{enumerate}
\item $\bTr^{(m,n)}$, from $\sig{\Sigma}^{(m)}$ to $\sig{\Sigma}^{(n)}$, with $\id_{S}$ on sorts, which sends every operation symbol to itself, applied to variables, and with domains $\Dom_{s}=\{(v^{s}_{0},\src^{m}_{s}(v^{s}_{0}))\}$; and
\item $\bTr^{\gl(m,n)}$, from $\sig{\Sigma}^{\gl(m)}$ to $\sig{\Sigma}^{\gl(n)}$, given by the canonical embedding, with empty domains.
\end{enumerate}
For $m=n$, both are identities.
The \emph{truncation functors} are $\Tr^{(m,n)}=(\bTr^{(m,n)})^{\ast}$ and $\Tr^{\gl(m,n)}=(\bTr^{\gl(m,n)})^{\ast}$.
\end{definition}

Thus the cells of $\Tr^{(m,n)}(\mathbf{B})$ are the $m$-cells of $\mathbf{B}$, and $\Tr^{\gl(m,n)}(\mathbf{B})$ is the reduct of $\mathbf{B}$ to the dimensions at most $m$.

\begin{proposition}\label{PTruncDerivor}
The functors $\Tr^{(m,n)}$ and $\Tr^{\gl(m,n)}$ restrict to functors from $\mathsf{nCatAlg}(\sig{\Sigma})$ to $\mathsf{mCatAlg}(\sig{\Sigma})$ and from $\mathsf{nCat}^{\gl}\mathsf{Alg}(\sig{\Sigma})$ to $\mathsf{mCat}^{\gl}\mathsf{Alg}(\sig{\Sigma})$.
\end{proposition}
\begin{proof}
For $m=n$ it is trivial.
Let $m<n$ and $\calg{B}$ an $n$-categorial $\sig{\Sigma}$-algebra.
By Lemma~\ref{LGradCells}, the $m$-cells are closed under the operations of $\Sigma^{(m)}$, and, by (A0), (A$^{k}$1) and (A$^{k}$3), these operations are defined on the families of $m$-cells that satisfy their domain conditions.
Hence, by Proposition~\ref{PTransformedEpAlg}, $\Tr^{(m,n)}(\calg{B})$ is the relative subalgebra of $\calg{B}$ on the $m$-cells, and it is a $\sig{\Sigma}^{(m)}$-algebra.
The axioms of $\mathcal{C}^{(m)}$ are axioms of $\mathcal{C}^{(n)}$, so that they hold in it.
In the globular case, the reduct is a $\sig{\Sigma}^{\gl(m)}$-algebra, since the domain conditions are the same, and the axioms of $\mathcal{C}^{\gl(m)}$ are axioms of $\mathcal{C}^{\gl(n)}$.
\end{proof}

\subsection{Discretization is left adjoint to truncation}

\begin{proposition}\label{PTruncDisc}
The equalities $\Tr^{(m,n)}\circ\Ds^{(n,m)}=\Id$ and $\Tr^{\gl(m,n)}\circ\Ds^{\gl(n,m)}=\Id$ hold.
\end{proposition}
\begin{proof}
Let $m<n$.
We first consider the single-set case.
The underlying derivor of $\bDs^{(n,m)}\circ\bTr^{(m,n)}$ is the identity, since both derivors send the operation symbols of $\Sigma^{(m)}$ to themselves.
Its domains are $(\Ds^{(n,m)})^{\sharp}[\{(v_{0},\src^{m}(v_{0}))\}]=\{(v_{0},v_{0})\}$, which every cell satisfies.
By Proposition~\ref{PDiscDerivor}, $\src^{m}$ is the identity in $\Ds^{(n,m)}(\calg{A})$, so that every cell of $\calg{A}$ is a cell of $\Tr^{(m,n)}(\Ds^{(n,m)}(\calg{A}))$.
By Propositions~\ref{PDiscDerivor} and~\ref{PTruncDerivor}, the latter is a $\sig{\Sigma}^{(m)}$-algebra.
Hence, by Proposition~\ref{PHepComp}~(2) and Lemma~\ref{LSameCells}, $\Tr^{(m,n)}(\Ds^{(n,m)}(\calg{A}))=(\id_{\sig{\Sigma}^{(m)}})^{\ast}(\calg{A})=\calg{A}$.
In the globular case, $\bDs^{\gl(n,m)}\circ\bTr^{\gl(m,n)}=\id_{\sig{\Sigma}^{\gl(m)}}$, since $\min\{k,m\}=k$ for $k\leq m$, and the equality follows in the same way.
\end{proof}

\begin{proposition}\label{PDiscTruncAdj}
For every $n$-categorial $\sig{\Sigma}$-algebra $\calg{B}$, let $\varepsilon_{\calg{B}}$ be the inclusion of the $m$-cells of $\calg{B}$, and, for every globular one, let $\varepsilon^{\gl}_{\calg{B}}$ be the identity on the dimensions $k\leq m$ and $\ii^{(k,m)}_{s}$ on the dimensions $k>m$.
Then $(\Tr^{(m,n)},\Ds^{(n,m)},\id,\varepsilon)$ and $(\Tr^{\gl(m,n)},\Ds^{\gl(n,m)},\id,\varepsilon^{\gl})$ are adjunctions.
Thus discretization is left adjoint to truncation.
\end{proposition}
\begin{proof}
Let $m<n$.
We first prove that $\varepsilon_{\calg{B}}$ is a homomorphism from $\Ds^{(n,m)}(\Tr^{(m,n)}(\calg{B}))$ to $\calg{B}$.
On $\Sigma$ and on the dimensions $k<m$, the domain is a relative subalgebra of $\calg{B}$.
For $k\geq m$ and an $m$-cell $x$, by Lemma~\ref{LGradCells}~(3) and (A$^{k}$5), $\src^{k}(x)=x=\tgt^{k}(x)$ and $x\circ^{k}x=x$ in $\calg{B}$, as in the domain.
Naturality is clear, since the functors act on homomorphisms by restriction.

In the globular case, $\varepsilon^{\gl}_{\calg{B}}$ is a homomorphism, by (EA$^{m,k}$7) for $\sigma_{k}$, by (EA$^{j,m,k}$1), (EA$^{j,m,k}$2), (EA$^{j,m,k}$5) and Lemma~\ref{LUnifInj} for $j<m$, and, for $j\geq m$, by (EA$^{m,j,k}$2), Lemma~\ref{LUnifInj} and (EA$^{j,k}$5), which give $\src^{(j,k)}(\ii^{(k,m)}(x))=\ii^{(j,m)}(x)$ and $\ii^{(k,m)}(x)\circ^{(j,k)}\ii^{(k,m)}(x)=\ii^{(k,m)}(x)$.
It is natural, since the homomorphisms commute with the identities.

Now we prove the triangular identities.
Since the unit is the identity, they state that $\Tr^{(m,n)}(\varepsilon_{\calg{B}})$ and $\varepsilon_{\Ds^{(n,m)}(\calg{A})}$ are identities.
The first one is the restriction of $\varepsilon_{\calg{B}}$ to the $m$-cells.
For the second one, every cell of $\Ds^{(n,m)}(\calg{A})$ is an $m$-cell.
In the globular case, the components of $\varepsilon^{\gl}_{\calg{B}}$ at the dimensions $k\leq m$ are identities, and the $(k,m)$-identities of $\Ds^{\gl(n,m)}(\calg{A})$ are identities.
\end{proof}

\subsection{Composites and compatibility with graduation and unification}

\begin{proposition}\label{PTruncComp}
Let $m\leq n\leq p$ in $\mathbb{N}$.
Then the following equalities hold:
\begin{enumerate}
\item $\Tr^{(m,n)}\circ\Tr^{(n,p)}=\Tr^{(m,p)}$ and $\Tr^{\gl(m,n)}\circ\Tr^{\gl(n,p)}=\Tr^{\gl(m,p)}$;
\item $\Ds^{(p,n)}\circ\Ds^{(n,m)}=\Ds^{(p,m)}$ and $\Ds^{\gl(p,n)}\circ\Ds^{\gl(n,m)}=\Ds^{\gl(p,m)}$; and
\item for every $q\leq p$ with $m\leq p$, $\Tr^{(q,p)}\circ\Ds^{(p,m)}$ is $\Ds^{(q,m)}$, if $m\leq q$, and $\Tr^{(q,m)}$, if $q\leq m$, and similarly in the globular case.
\end{enumerate}
\end{proposition}
\begin{proof}
In every case we compare the composite generalized hep-derivors and then apply Proposition~\ref{PHepComp}~(2) and Lemma~\ref{LSameCells}.

(1) The underlying derivors agree, since all of them are inclusions.
In the single-set case, the domain of the composite is $\{(v_{0},\src^{n}(v_{0})),(v_{0},\src^{m}(v_{0}))\}$, which determines the $m$-cells, by Lemma~\ref{LGradCells}~(3).
In the globular case, all the domains are empty.

(2) The composite of the underlying derivors sends a symbol of dimension $k$ to itself, if $k<m$, and to a variable, otherwise; and $\min\{\min\{k,n\},m\}=\min\{k,m\}$.
All the domains are empty.

(3) If $m\leq q$, the composite derivor is $\bDs^{(q,m)}$ on terms; if $q\leq m$, it is $\bTr^{(q,m)}$ on terms.
In both cases the cells of $\Tr^{(q,p)}(\Ds^{(p,m)}(\calg{A}))$ are, by Proposition~\ref{PDiscDerivor}, all the cells of $\calg{A}$, if $m\leq q$, and its $q$-cells, if $q\leq m$.
The globular case is similar.
\end{proof}

\begin{proposition}\label{PCompatGdUf}
The following equalities hold:
\begin{enumerate}
\item $\Tr^{\gl(m,n)}\circ\Gd^{(n)}=\Gd^{(m)}\circ\Tr^{(m,n)}$;
\item $\Gd^{(n)}\circ\Ds^{(n,m)}=\Ds^{\gl(n,m)}\circ\Gd^{(m)}$; and
\item $\Uf^{(n)}\circ\Ds^{\gl(n,m)}=\Ds^{(n,m)}\circ\Uf^{(m)}$.
\end{enumerate}
Moreover, for every globular $n$-categorial $\sig{\Sigma}$-algebra $\calg{B}$, the family $\beta^{(m,n)}_{\calg{B}}=(\ii^{(n,m)}_{s})_{s\in S}$ is a natural isomorphism from $\Uf^{(m)}(\Tr^{\gl(m,n)}(\calg{B}))$ to $\Tr^{(m,n)}(\Uf^{(n)}(\calg{B}))$.
\end{proposition}
\begin{proof}
All the functors involved act on homomorphisms by restriction, so that it suffices to compare the values on objects.

(1) By Proposition~\ref{PGradDerivor} and Lemma~\ref{LGradCells}~(3), in both algebras the cells of sort $(k,s)$ are the $k$-cells of sort $s$, for $k\leq m$, with the restricted operations.

(2) By Propositions~\ref{PGradDerivor} and~\ref{PDiscDerivor}, in both algebras the cells of sort $(k,s)$ are the $\min\{k,m\}$-cells of $\calg{A}$, with the structure of $\calg{A}$ for $j<m$ and the trivial structure for $j\geq m$.

(3) By Propositions~\ref{PUnifDerivor} and~\ref{PDiscDerivor}, in both algebras the cells are the $m$-cells of $\calg{B}$, with $\src^{k}=\ii^{(m,k)}\circ\src^{(k,m)}$ and $\circ^{k}=\circ^{(k,m)}$ for $k<m$, and the trivial structure for $k\geq m$.

Finally, the cells of $\Tr^{(m,n)}(\Uf^{(n)}(\calg{B}))$ are the $n$-cells $x$ with $x=\ii^{(n,m)}(\src^{(m,n)}(x))$, i.e., by Lemma~\ref{LUnifInj}, the image of $\ii^{(n,m)}$, which is injective.
By (EA$^{m,n}$7), $\beta^{(m,n)}_{\calg{B}}$ commutes with the operations of $\Sigma$.
By (EA$^{k,m,n}$1), (EA$^{k,m,n}$2) and Lemma~\ref{LUnifInj}, it commutes with the sources and the targets.
By (EA$^{k,m,n}$5), it preserves and reflects the compositions.
Hence it is an isomorphism, natural since the homomorphisms commute with the identities.
\end{proof}

\section{The dimension \texorpdfstring{$\omega$}{omega}}\label{SOmega}

In this section we extend the above functors to the dimension $\omega$.
Truncation and discretization are again induced by generalized hep-derivors, and they present the $\omega$-categorial algebras as projective limits.
Graduation is then obtained from the universal property of the limit, and unification as an inductive limit.
Let $\sig{\Sigma}=(S,\Sigma)$ be a many-sorted signature and $n\in\mathbb{N}$.

\subsection{Truncation and discretization}

We let $\bTr^{(n,\omega)}$, $\bTr^{\gl(n,\omega)}$, $\bDs^{(\omega,n)}$ and $\bDs^{\gl(\omega,n)}$ stand for the generalized hep-derivors defined by the clauses of Definitions~\ref{DDiscDerivor} and~\ref{DTruncDerivor}, with $\omega$ in place of the greatest dimension.
For every $p>n$, the restriction of $\bDs^{(\omega,n)}$ to $\Sigma^{(p)}$ is $\bDs^{(p,n)}$, and similarly in the globular case.
Hence, by Proposition~\ref{PTransformedAlg}, the reduct of $\Ds^{(\omega,n)}(\calg{B})$ to $\Sigma^{(p)}$ is $\Ds^{(p,n)}(\calg{B})$.

\begin{proposition}\label{PTruncOmega}
The following functors are well defined:
\begin{enumerate}
\item $\Tr^{(n,\omega)}\colon\omega\mathsf{CatAlg}(\sig{\Sigma})\mor\mathsf{nCatAlg}(\sig{\Sigma})$ and $\Tr^{\gl(n,\omega)}\colon\omega\mathsf{Cat}^{\gl}\mathsf{Alg}(\sig{\Sigma})\mor\mathsf{nCat}^{\gl}\mathsf{Alg}(\sig{\Sigma})$; and
\item $\Ds^{(\omega,n)}\colon\mathsf{nCatAlg}(\sig{\Sigma})\mor\omega\mathsf{CatAlg}(\sig{\Sigma})$ and $\Ds^{\gl(\omega,n)}\colon\mathsf{nCat}^{\gl}\mathsf{Alg}(\sig{\Sigma})\mor\omega\mathsf{Cat}^{\gl}\mathsf{Alg}(\sig{\Sigma})$.
\end{enumerate}
Moreover, the equalities of Proposition~\ref{PTruncComp} and the adjunctions of Proposition~\ref{PDiscTruncAdj} hold with $\omega$ in place of $p$ and $n$, respectively.
\end{proposition}
\begin{proof}
(1) The proof of Proposition~\ref{PTruncDerivor} only uses Lemma~\ref{LGradCells}, which holds in dimension $\omega$, and axioms for finitely many dimensions.

(2) We first prove that $\Ds^{(\omega,n)}(\calg{B})$ is a $\sig{\Sigma}^{(\omega)}$-algebra in which the axioms hold.
Every operation symbol, and every axiom, involves finitely many dimensions.
Choose $p>n$ greater than all of them.
Then the operation symbol, or the axiom, is interpreted as in $\Ds^{(p,n)}(\calg{B})$, which, by Proposition~\ref{PDiscDerivor}, is a $p$-categorial algebra.
Now we prove that every cell is an $n$-cell.
In fact, by Proposition~\ref{PDiscDerivor}, $\src^{n}$ is the identity.
The globular case is the same.

The equalities of Proposition~\ref{PTruncComp} only involve finitely many dimensions at a time, and their proofs carry over.
For the adjunction, the counit at an $\omega$-categorial algebra is again the inclusion of the $m$-cells, or the family of identities and $(k,m)$-identities, and the proof of Proposition~\ref{PDiscTruncAdj} carries over.
\end{proof}

\subsection{Projective limits}

By Proposition~\ref{PTruncOmega}, the categories $\mathsf{nCatAlg}(\sig{\Sigma})$, with the functors $\Tr^{(m,n)}$, form a projective system of categories, and the functors $\Tr^{(n,\omega)}$ form a cone over it; and similarly in the globular case.

\begin{proposition}\label{PUnivPropCatAlg}
Let $\mathsf{L}$ be a category and $(F_{n}\colon\mathsf{L}\mor\mathsf{nCatAlg}(\sig{\Sigma}))_{n\in\omega}$ a family of functors with $F_{m}=\Tr^{(m,n)}\circ F_{n}$ for $m\leq n$.
Then there exists a unique functor $\langle F_{k}\rangle_{k}\colon\mathsf{L}\mor\omega\mathsf{CatAlg}(\sig{\Sigma})$ with $\Tr^{(n,\omega)}\circ\langle F_{k}\rangle_{k}=F_{n}$ for every $n$.
The same holds in the globular case.
Hence $\omega\mathsf{CatAlg}(\sig{\Sigma})$ and $\omega\mathsf{Cat}^{\gl}\mathsf{Alg}(\sig{\Sigma})$ are the projective limits of the categories of $n$-categorial algebras.
\end{proposition}
\begin{proof}
We first consider the single-set case.
Let $(\calg{A}_{n})_{n}$ be a compatible family of objects, i.e., $\Tr^{(m,n)}(\calg{A}_{n})=\calg{A}_{m}$ for $m\leq n$.
Then the cells of each sort form an increasing chain.
We let $\varprojlim_{n}\calg{A}_{n}$ be the union, in which every operation, on finitely many cells and of dimension $k$, is computed in any $\calg{A}_{n}$ with $n>k$ that contains the cells.
This does not depend on $n$, since $\calg{A}_{n}$ is the relative subalgebra of $\calg{A}_{n'}$ on its $n$-cells.
Every axiom involves finitely many cells and dimensions, so that it holds.
Every cell of $\calg{A}_{n}$ is an $n$-cell of the limit, since $\calg{A}_{n}=\Tr^{(n,n+1)}(\calg{A}_{n+1})$.
Hence the limit is an $\omega$-categorial algebra, and its $n$-cells are the cells of $\calg{A}_{n}$, i.e., $\Tr^{(n,\omega)}(\varprojlim_{k}\calg{A}_{k})=\calg{A}_{n}$.
Conversely, every $\omega$-categorial algebra is the union of its truncations, so that it is determined by them.
For a compatible family of homomorphisms $(H_{n})_{n}$, the mapping that acts on the cells of $\calg{A}_{n}$ as $H_{n}$ is the unique homomorphism with truncations $H_{n}$.
Now let $\langle F_{k}\rangle_{k}(X)=\varprojlim_{n}F_{n}(X)$ and $\langle F_{k}\rangle_{k}(f)=\varprojlim_{n}F_{n}(f)$.
By the uniqueness of the limits, it is a functor, it satisfies $\Tr^{(n,\omega)}\circ\langle F_{k}\rangle_{k}=F_{n}$, and it is the unique one.

In the globular case, the limit of a compatible family $(\calg{A}_{n})_{n}$ has as $k$-cells the $k$-cells of $\calg{A}_{k}$, with the operations of the $\calg{A}_{k}$, and the same argument applies.
\end{proof}

\begin{corollary}\label{CLimitNat}
Let $F,G\colon\mathsf{L}\mor\omega\mathsf{CatAlg}(\sig{\Sigma})$ be functors and $(\varphi^{n}\colon\Tr^{(n,\omega)}\circ F\mor\Tr^{(n,\omega)}\circ G)_{n}$ natural transformations with $\Tr^{(m,n)}\ast\varphi^{n}=\varphi^{m}$.
Then there exists a unique natural transformation $\varphi\colon F\mor G$ with $\Tr^{(n,\omega)}\ast\varphi=\varphi^{n}$, and it is a natural isomorphism if every $\varphi^{n}$ is.
The same holds in the globular case.
\end{corollary}
\begin{proof}
By the proof of Proposition~\ref{PUnivPropCatAlg}, every compatible family of homomorphisms has a unique limit.
The limit of $(\varphi^{n}_{X})_{n}$ is natural in $X$, since $G(f)\circ\varphi_{X}$ and $\varphi_{Y}\circ F(f)$ have the same truncations.
If every $\varphi^{n}_{X}$ is invertible, the limit of the inverses is the inverse.
\end{proof}

\subsection{Graduation and unification}

\begin{definition}\label{DGdOmega}
The \emph{graduation functor} $\Gd^{(\omega)}\colon\omega\mathsf{CatAlg}(\sig{\Sigma})\mor\omega\mathsf{Cat}^{\gl}\mathsf{Alg}(\sig{\Sigma})$ is the unique functor such that $\Tr^{\gl(n,\omega)}\circ\Gd^{(\omega)}=\Gd^{(n)}\circ\Tr^{(n,\omega)}$, for every $n\in\mathbb{N}$.
\end{definition}

It exists by Proposition~\ref{PUnivPropCatAlg}, since, by Propositions~\ref{PCompatGdUf} and~\ref{PTruncOmega}, the family $(\Gd^{(n)}\circ\Tr^{(n,\omega)})_{n}$ is compatible.
Explicitly, the $k$-cells of $\Gd^{(\omega)}(\calg{A})$ are the $k$-cells of $\calg{A}$, with the operations of Proposition~\ref{PGradDerivor}.

Unification cannot be obtained in the same way.
It is not induced by a generalized hep-derivor: the cells of a sort of the transformed algebra are cells of a single sort $(k,s')$, whereas a globular $\omega$-category with one cell in every dimension $j\leq k$ and two $(k+1)$-cells has a unification with two cells.
Nor is the family $(\Uf^{(n)}\circ\Tr^{\gl(n,\omega)})_{n}$ compatible: by Proposition~\ref{PCompatGdUf}, it is only compatible up to the isomorphisms $\beta^{(m,n)}$, which are, however, the inclusions of an inductive system.

\begin{definition}\label{DUfOmega}
Let $\calg{B}$ be a globular $\omega$-categorial $\sig{\Sigma}$-algebra and $s\in S$.
We let $[B]_{s}$ stand for the quotient of $\coprod_{k\in\omega}B_{(k,s)}$ by the equivalence relation that identifies $x\in B_{(j,s)}$ with $\ii^{(k,j)}_{s}(x)\in B_{(k,s)}$, for $j<k$, and $[x]_{k}$ for the class of $x\in B_{(k,s)}$.
The \emph{unification} $\Uf^{(\omega)}(\calg{B})$ is the partial $\Sigma^{(\omega)}$-algebra on $[B]$ in which every operation symbol $\tau$ of $\Sigma^{(p)}$ acts on classes $[x_{j}]_{p}$ represented at the level $p$ as $\tau$ acts on the $x_{j}$ in $\Uf^{(p)}(\Tr^{\gl(p,\omega)}(\calg{B}))$, with the same domain.
For a homomorphism $F$, $\Uf^{(\omega)}(F)$ sends $[x]_{k}$ to $[F_{(k,s)}(x)]_{k}$.
\end{definition}

By (EA$^{h,j,k}$2), the above relation is an equivalence, and, by Lemma~\ref{LUnifInj}, every mapping $x\mapsto[x]_{k}$, denoted by $\mathrm{in}^{[k]}_{s}$, is injective.
Every finite family of classes is represented at a common level.
By Proposition~\ref{PCompatGdUf}, $\ii^{(p',p)}$ is an isomorphism from the level $p$ onto the $p$-cells of the level $p'$, so that the definition does not depend on the level.

\begin{proposition}\label{PUfOmega}
For every globular $\omega$-categorial $\sig{\Sigma}$-algebra $\calg{B}$, $\Uf^{(\omega)}(\calg{B})$ is an $\omega$-categorial $\sig{\Sigma}$-algebra whose $k$-cells are the classes $[x]_{k}$, and $\mathrm{in}^{[n]}$ is a natural isomorphism from $\Uf^{(n)}\circ\Tr^{\gl(n,\omega)}$ to $\Tr^{(n,\omega)}\circ\Uf^{(\omega)}$.
Moreover, $\Uf^{(\omega)}$ is a functor.
\end{proposition}
\begin{proof}
We first prove that the axioms of $\mathcal{C}^{(\omega)}$ hold.
Every axiom involves finitely many dimensions and variables.
For a valuation, we choose a level $p$ greater than these dimensions at which the values are represented.
Then the axiom holds, since it holds in the $p$-categorial algebra $\Uf^{(p)}(\Tr^{\gl(p,\omega)}(\calg{B}))$.

Now we compute the $k$-cells.
For $x\in B_{(k,s)}$ and $p>k$, by Proposition~\ref{PUnifDerivor} and Lemma~\ref{LUnifInj}, $\src^{k}(\ii^{(p,k)}(x))=\ii^{(p,k)}(\src^{(k,p)}(\ii^{(p,k)}(x)))=\ii^{(p,k)}(x)$, so that $[x]_{k}$ is a $k$-cell.
Conversely, if $[y]_{p}$ is a $k$-cell, then $[y]_{p}=[\src^{(k,p)}(y)]_{k}$.
Hence every cell has a finite dimension.

Next, $\mathrm{in}^{[n]}$ is a bijection onto the $n$-cells, and, by Definition~\ref{DUfOmega}, it preserves and reflects the operations of $\Sigma^{(n)}$.
Finally, $\Uf^{(\omega)}(F)\circ\mathrm{in}^{[p]}=\mathrm{in}^{[p]}\circ\Uf^{(p)}(\Tr^{\gl(p,\omega)}(F))$, so that $\Uf^{(\omega)}(F)$ is a homomorphism, and $\Uf^{(\omega)}$ preserves identities and composites.
\end{proof}

\begin{proposition}\label{POmegaGdUf}
The functors $\Gd^{(\omega)}$ and $\Uf^{(\omega)}$ are mutually quasi-inverse equivalences, and there exists an adjoint equivalence in which $\Gd^{(\omega)}$ is left adjoint to $\Uf^{(\omega)}$.
\end{proposition}
\begin{proof}
We first prove that $\Uf^{(\omega)}\circ\Gd^{(\omega)}\cong\Id$.
Let $\calg{A}$ be an $\omega$-categorial algebra.
The identities of $\Gd^{(\omega)}(\calg{A})$ are inclusions, so that $[x]_{j}=[y]_{k}$ if, and only if, $x=y$.
Hence $\zeta_{\calg{A}}\colon[x]_{k}\mapsto x$ is a bijection onto the cells of $\calg{A}$.
For every level $p$, by Definition~\ref{DGdOmega} and Proposition~\ref{PUnifGrad}, $\Uf^{(p)}(\Tr^{\gl(p,\omega)}(\Gd^{(\omega)}(\calg{A})))=\Tr^{(p,\omega)}(\calg{A})$.
Hence $\zeta_{\calg{A}}$ preserves and reflects the operations, and it is a natural isomorphism.

Now we prove that $\Gd^{(\omega)}\circ\Uf^{(\omega)}\cong\Id$.
Let $\calg{B}$ be a globular $\omega$-categorial algebra.
By Proposition~\ref{PUfOmega}, the family $(\mathrm{in}^{[k]}_{s})_{(k,s)}$ is a bijection from the cells of $\calg{B}$ onto those of $\Gd^{(\omega)}(\Uf^{(\omega)}(\calg{B}))$.
It commutes with the operations, since the $j$-source of $[x]_{k}$ is $[\src^{(j,k)}(x)]_{j}$, the $(k,j)$-identity of $[u]_{j}$ is $[u]_{j}=[\ii^{(k,j)}(u)]_{k}$, and the compositions and the $\sigma$ are computed on representatives.
Hence it is a natural isomorphism.
Finally, every equivalence can be completed to an adjoint equivalence, see~\cite{sM98}, Theorem~IV.4.1.
\end{proof}

\section{Algebraic squares}\label{SSquares}

Four adjunctions form an \emph{algebraic square}, see~\cite{CVCLRM26}, when they are arranged as follows and there exist natural isomorphisms between the composites of the left adjoints and between the composites of the right adjoints.

\begin{definition}\label{DAlgSquare}
Let $F\dashv G$, from $\mathsf{C}$ to $\mathsf{D}$, $F'\dashv G'$, from $\mathsf{C}'$ to $\mathsf{D}'$, $J\dashv K$, from $\mathsf{C}$ to $\mathsf{C}'$, and $H\dashv I$, from $\mathsf{D}$ to $\mathsf{D}'$, be adjunctions.
They form an \emph{algebraic square} if there exist natural isomorphisms $\alpha\colon F'\circ J\mor H\circ F$ and $\beta\colon G\circ I\mor K\circ G'$.
\end{definition}

\begin{proposition}\label{PAlgSquare}
Let $\sig{\Sigma}$ be a many-sorted signature and $m\leq n$ in $\mathbb{N}$.
Then the following statements hold:
\begin{enumerate}
\item the adjunctions $\Gd^{(m)}\dashv\Uf^{(m)}$, $\Gd^{(n)}\dashv\Uf^{(n)}$, $\Ds^{(n,m)}\dashv\Tr^{(m,n)}$ and $\Ds^{\gl(n,m)}\dashv\Tr^{\gl(m,n)}$ form an algebraic square, with $\alpha$ the identity and $\beta=\beta^{(m,n)}$; and
\item the same holds with $\omega$ in place of $n$, with $\alpha$ the identity and $\beta=\mathrm{in}^{[m]}$.
\end{enumerate}
\end{proposition}
\begin{proof}
(1) The adjunctions are those of Propositions~\ref{PGdUfAdj} and~\ref{PDiscTruncAdj}.
By Proposition~\ref{PCompatGdUf}, $\Gd^{(n)}\circ\Ds^{(n,m)}=\Ds^{\gl(n,m)}\circ\Gd^{(m)}$ and $\beta^{(m,n)}$ is a natural isomorphism from $\Uf^{(m)}\circ\Tr^{\gl(m,n)}$ to $\Tr^{(m,n)}\circ\Uf^{(n)}$.

(2) The adjunctions are those of Propositions~\ref{PGdUfAdj}, \ref{POmegaGdUf} and~\ref{PTruncOmega}.
For every $p\in\mathbb{N}$, by Definition~\ref{DGdOmega} and Propositions~\ref{PCompatGdUf} and~\ref{PTruncOmega}, $\Gd^{(\omega)}\circ\Ds^{(\omega,m)}$ and $\Ds^{\gl(\omega,m)}\circ\Gd^{(m)}$ have the same truncations.
Hence, by Proposition~\ref{PUnivPropCatAlg}, they are equal.
By Proposition~\ref{PUfOmega}, $\mathrm{in}^{[m]}$ is a natural isomorphism from $\Uf^{(m)}\circ\Tr^{\gl(m,\omega)}$ to $\Tr^{(m,\omega)}\circ\Uf^{(\omega)}$.
\end{proof}

\section{Change of signature}\label{SChangeSig}

In this section we prove that all the above constructions are compatible with the derivors between many-sorted signatures, regarded as hep-signatures with empty domain conditions.
Let $\sig{\Sigma}=(S,\Sigma)$ and $\sig{\Lambda}=(T,\Lambda)$ be many-sorted signatures, $\mathbf{c}=(\varphi,c)$ a derivor from $\sig{\Sigma}$ to $\sig{\Lambda}$ and $n\in\mathbb{N}$.
For a term $P$ of $\Lambda$ and $k\in n+1$, we let $P_{k}$ be as in Lemma~\ref{LTermFunctor}.

\begin{definition}\label{DChangeSigDerivor}
The generalized hep-derivors determined by $\mathbf{c}$ are the following ones, with empty domains:
\begin{enumerate}
\item $\mathbf{c}^{(n)}$, from $\sig{\Sigma}^{(n)}$ to $\sig{\Lambda}^{(n)}$, with $\varphi$ on sorts, which sends $\sigma$ to $c(\sigma)$, and $\src^{k}_{s}$, $\tgt^{k}_{s}$ and $\circ^{k}_{s}$ to $\src^{k}_{\varphi(s)}$, $\tgt^{k}_{\varphi(s)}$ and $\circ^{k}_{\varphi(s)}$, applied to variables; and
\item $\mathbf{c}^{\gl(n)}$, from $\sig{\Sigma}^{\gl(n)}$ to $\sig{\Lambda}^{\gl(n)}$, which sends $(k,s)$ to $(k,\varphi(s))$, $\sigma_{k}$ to $c(\sigma)_{k}$, and $\src^{(j,k)}_{s}$, $\tgt^{(j,k)}_{s}$, $\ii^{(k,j)}_{s}$ and $\circ^{(j,k)}_{s}$ to the same symbols of sort $\varphi(s)$, applied to variables.
\end{enumerate}
For $n=0$, $\mathbf{c}^{(0)}$ is $\mathbf{c}$.
The same clauses, with $\omega$ in place of $n$, define $\mathbf{c}^{(\omega)}$ and $\mathbf{c}^{\gl(\omega)}$.
\end{definition}

\begin{proposition}\label{PChangeSigDerivor}
The functors $(\mathbf{c}^{(n)})^{\ast}$ and $(\mathbf{c}^{\gl(n)})^{\ast}$ restrict to functors from $\mathsf{nCatAlg}(\sig{\Lambda})$ to $\mathsf{nCatAlg}(\sig{\Sigma})$ and from $\mathsf{nCat}^{\gl}\mathsf{Alg}(\sig{\Lambda})$ to $\mathsf{nCat}^{\gl}\mathsf{Alg}(\sig{\Sigma})$.
The first one keeps, in every sort $s$, the $n$-category of sort $\varphi(s)$, and interprets $\sigma$ as the $n$-functor $c(\sigma)^{\calg{B}}$; the second one is similar.
Moreover, $(\mathbf{d}\circ\mathbf{c})^{(n)}=\mathbf{d}^{(n)}\circ\mathbf{c}^{(n)}$, $(\id_{\sig{\Sigma}})^{(n)}=\id_{\sig{\Sigma}^{(n)}}$, and similarly in the globular case, so that $\sig{\Sigma}\mapsto\mathsf{nCatAlg}(\sig{\Sigma})$ and $\sig{\Sigma}\mapsto\mathsf{nCat}^{\gl}\mathsf{Alg}(\sig{\Sigma})$ are split $\mathsf{Sig}_{\mathfrak{d}}$-indexed categories.
The same holds in dimension $\omega$.
\end{proposition}
\begin{proof}
Let $\calg{B}$ be an $n$-categorial $\sig{\Lambda}$-algebra.
We first prove that $(\mathbf{c}^{(n)})^{\ast}(\calg{B})$ is a $\sig{\Sigma}^{(n)}$-algebra.
The domains are empty, so that its cells of sort $s$ are the elements of $B_{\varphi(s)}$.
By Lemma~\ref{LTermFunctor}, $c(\sigma)^{\calg{B}}$ is total, and so are the sources and the targets.
By (A$^{k}$3), the composition is defined on the pairs that satisfy its domain condition.

Now we prove the axioms.
The axioms on the categorial structure of the sort $s$ are those of the sort $\varphi(s)$ of $\calg{B}$.
(A0), (A$^{k}$7) and (A$^{k}$8) state that $c(\sigma)^{\calg{B}}$ is a total $n$-functor, which is Lemma~\ref{LTermFunctor}.
The globular case is the same.

Next, we prove the functoriality.
We compare the composite derivors on operation symbols, see the note after Proposition~\ref{PDerivors}.
For every term $P$ of $\Lambda$, $(d^{(n)})^{\sharp}(P)=d^{\sharp}(P)$ and $(d^{\gl(n)})^{\sharp}(P_{k})=(d^{\sharp}(P))_{k}$, since both sides of each equality are composites of derivors that agree on the operation symbols of $\Lambda$.
Hence both composites send $\sigma$ to $d^{\sharp}(c(\sigma))$, and $\sigma_{k}$ to $(d^{\sharp}(c(\sigma)))_{k}$, and the categorial symbols to themselves.
The identities are clear.
By Proposition~\ref{PHepComp}~(2), since all the domains are empty and the transformed algebras are ep-algebras, the induced functors compose strictly.

In dimension $\omega$, the proofs only involve finitely many dimensions at a time.
Moreover, the cells of $(\mathbf{c}^{(\omega)})^{\ast}(\calg{B})$ have the dimensions they have in $\calg{B}$, so that $(\mathbf{c}^{(\omega)})^{\ast}$ preserves the finite dimension of the cells.
\end{proof}

\begin{proposition}\label{PChangeSigCompat}
The functors induced by $\mathbf{c}$ commute with graduation, unification, discretization and truncation:
\begin{enumerate}
\item $\Gd^{(n)}_{\sig{\Sigma}}\circ(\mathbf{c}^{(n)})^{\ast}=(\mathbf{c}^{\gl(n)})^{\ast}\circ\Gd^{(n)}_{\sig{\Lambda}}$ and $\Uf^{(n)}_{\sig{\Sigma}}\circ(\mathbf{c}^{\gl(n)})^{\ast}=(\mathbf{c}^{(n)})^{\ast}\circ\Uf^{(n)}_{\sig{\Lambda}}$;
\item $\Ds^{(n,m)}_{\sig{\Sigma}}\circ(\mathbf{c}^{(m)})^{\ast}=(\mathbf{c}^{(n)})^{\ast}\circ\Ds^{(n,m)}_{\sig{\Lambda}}$ and $\Tr^{(m,n)}_{\sig{\Sigma}}\circ(\mathbf{c}^{(n)})^{\ast}=(\mathbf{c}^{(m)})^{\ast}\circ\Tr^{(m,n)}_{\sig{\Lambda}}$, for $m\leq n$, and similarly in the globular case; and
\item the same equalities hold in dimension $\omega$.
\end{enumerate}
Moreover, the natural transformations $\theta$, $\beta^{(m,n)}$, $\varepsilon$, $\varepsilon^{\gl}$ and $\mathrm{in}^{[m]}$ are compatible with the change of signature.
\end{proposition}
\begin{proof}
(1) We first prove that $\mathbf{c}^{(n)}\circ\bGd^{(n)}_{\sig{\Sigma}}=\bGd^{(n)}_{\sig{\Lambda}}\circ\mathbf{c}^{\gl(n)}$.
Both send $(k,s)$ to $\varphi(s)$.
By (Der1), (Der2) and the fact that $(\Gd^{(n)}_{\sig{\Lambda}})^{\sharp}(P_{k})=P$, both send $\sigma_{k}$ to $c(\sigma)$; and both send every categorial symbol to the corresponding term of Definition~\ref{DGradDerivor}, for the sort $\varphi(s)$.
Both have the domains $\{(v_{0},\src^{k}_{\varphi(s)}(v_{0}))\}$, for $k\in n$.
Moreover, both $\Gd^{(n)}_{\sig{\Sigma}}((\mathbf{c}^{(n)})^{\ast}(\calg{B}))$ and $(\mathbf{c}^{\gl(n)})^{\ast}(\Gd^{(n)}_{\sig{\Lambda}}(\calg{B}))$ are ep-algebras whose cells of sort $(k,s)$ are the $k$-cells of sort $\varphi(s)$ of $\calg{B}$.
Hence, by Lemma~\ref{LCommSquare}, they are equal.
For unification, both derivors send $s$ to $(n,\varphi(s))$, $\sigma$ to $c(\sigma)_{n}$, since $(\Uf^{(n)}_{\sig{\Lambda}})^{\sharp}(P)=P_{n}$, and the categorial symbols to the terms of Definition~\ref{DUnifDerivor}; all the domains are empty, and the cells are the $n$-cells of sort $\varphi(s)$.

(2) In the same way, the composite derivors agree, since discretization and truncation fix the terms of $\Lambda$, up to replacing $P_{k}$ by $P_{\min\{k,m\}}$, and since both composites have the same domains.
The cells are, in both cases, the cells, the $m$-cells or the $\min\{k,m\}$-cells of the sort $\varphi(s)$.

(3) The equalities for $\Tr^{(n,\omega)}$ and $\Ds^{(\omega,n)}$ are proved in the same way.
For $\Gd^{(\omega)}$, both functors have the same truncations, by (1) and (2), so that, by Proposition~\ref{PUnivPropCatAlg}, they are equal.
For $\Uf^{(\omega)}$, both algebras have the classes $[x]_{k}$, with $x$ of sort $(k,\varphi(s))$, as cells of sort $s$, and their operations are computed at a finite level, where they agree by~(1).

Finally, these natural transformations are given by identities, inclusions, $(k,j)$-identities and the mappings $\mathrm{in}^{[m]}$, which the induced functors reindex along $\varphi$.
\end{proof}

Hence graduation, unification, discretization and truncation, and the algebraic squares of Proposition~\ref{PAlgSquare}, are natural in the many-sorted signature.

\bibliographystyle{amsplain}
\bibliography{Biblio}

\end{document}
